% ================================================================
%  18.5.7 -- 18.5.9 (原书 p.864–893)
%  源文: _work/chunks/chunk10.md, chunk11.md, chunk121.md, chunk122.md,
%        chunk123.md, chunk124.md
% ================================================================
\subsection{升级(2)}   % 18.5.7. L'escalade (2)

\notehead{174}{}

\dmark{}\origpage{1148}(3月22日)\fnote{889}{(*)}{(6月14日)眼下这道注,接续前一道注(同日所写)的 a.部分(“一位备受簇拥的死者”)而来。}我在别处已经强调过:其实并没有四种运算(对\textbf{一场}葬礼而言),只有
独一无二的一场“\textbf{葬礼运算}”。把它分成四大块,不过是为了叙述方便,其实是人为
的,而且(若太按字面当真)容易误导。确凿无疑,在那位导演——乐队指挥——葬礼主祭的身
体里,从来没有\textbf{四}只小魔鬼蹲在脑袋的四个角落里给他递台词,只有独一个!在关于阴与阳
的那场漫长冥想里\fnote{890}{(*)}{这就是构成《收获与播种》第三部分主体、Notes 104--162$''$的那场思考。},我曾试着与这只小魔鬼比从前更多地相识——从前,我不过时时察觉
它还在那儿折腾,转眼又去管别的。不过我并不声称已经充分认识了它;说到底,认识它或许根本
不是我的分内之事。确定无疑的是,它还像从前一样折腾个不停,而且未必会在我朋友咽气之前收
场。话虽如此,那场著名的“葬礼运算”至今仍在进行——就在我写这几行的此刻仍在进行。我不知
道,眼下这本“家庭相册”的散发,能不能至少终止全部局部运算中最庞大(也最不义)的那一项:
活埋一位年轻数学家——Zoghman~Mebkhout,“所有”在代数簇或复簇上同调里做工的人,四五年来
一直在用他的思想和结果……

既然分明只有一场运算,我就放下“四种”运算的虚构,不妨按时间先后,把我知道的主要插曲
与阶段勾一个轮廓。这里我不做了——我自认为已经做够了:在前四条主注(“沉默”、“操演”、
“分赃”、“封神”,n\textsuperscript{o}s~168、169、170、171)里,我把所知的全部插曲都收拢在一处,好奇
的读者自会把它在时间的标尺上排定次序。说来奇怪,若从“第二层次”或“运算”的观点(容我
用用这些委婉语)看,我1970年告别数学舞台的那一年,在这串插曲的接续之中竟不构成什么断
裂:在我印象里,它们自1966年 SGA~5 讨论班结束起,一直以相当均匀的节奏延续到1977年——
那一年,“SGA~4½--SGA~5”双双出版,SGA~5 又出了Illusie的\dmark{}\origpage{1149}编订版\fnote{891}{(*)}{(6月3日)这个印象应当修正,须把那场声势浩大的“Tannaka 范畴”(原文如此!)运算考虑进来:其第一插曲(与“稻草爸爸”N.~Saavedra搭档)在1972年(尾声则在1982年,由“真爸爸”Pierre~Deligne接力)。关于这一点,见那组注“第六根钉子(钉进棺材)”(n\textsuperscript{o}s~176$_{1}$--176$_{7}$)。}。这场运算,我看作一个骤然而惊人的\textbf{质变}。之前,是一次不动
声色的“收割”。从这里,我感到一阵暴烈与蔑视突然发作,死死缠住一个已被宣布“亡故”的缺席
者的作品。

在我那班上同调学生的这场上上下下齐心的集体发泄之后(“教团全体”在旁边看得心满意足),
似乎有过四年的平静。而从1966到1977的十一年间,我隔一两年就能辨认出一个类型分明的“插
曲”;1977到1981(“反常层会议”之年)之间,我却一无所知。相反,Deligne的长文《父亲猜想,
II》(La conjecture de père,~II),1980年发表于《数学公刊》——就在那场令人难以置信的讨论
会的前一年——照眼下的世道,几乎可以算作正常了……\fnote{892}{(*)}{当然,这项工作的主要结果没有一字提到我这个人;而那个结果的表述,本就是Deligne从我这里拿去的 Motif 瑜伽的一部分。反倒有一件事让我吃了一惊:我的名字,与Miller的名字一道,出现在文中一个段落里,说起那个除幂德拉姆复形——它大约在1976年由Miller和我各自独立地引入。1976年我曾在法国高等科学研究所就这个题目做过一次报告(顺便说,那是我一生中最后一次公开演讲),但明摆着我什么也不会发表。换了旁人,把这份纯属非正式的共同父亲身份略去不表,想必谁也不会察觉,更不会有人有半句闲话……}。也是在这一年,Deligne在一次Bourbaki讨论班上,从作者本
人嘴里得知了“好上帝定理”(又名Mebkhout定理)\fnote{893}{(*)}{(6月1日)其实,这一幕发生在前一年,即1979年6月,就在Bourbaki讨论班上。}。这是长期停滞的上同调主题冰雪骤消的开端。也是信号:第二年,葬礼
运算的第二场也是最后一场(?)高潮就此发动,这一回的调门是不义的,一切克制,连同最起码
的审慎,都被轻快地抛到了九霄云外。

翌年“难忘文集”LN~900那一幕(它为 Motif 掘墓复生,只字不提我这个人——去年某个4月19
日,这一幕曾让我百感交集……),还有同年Berthelot的报告(它把卑不足道的我从晶体上同调的
“历史”——原文如此——中一举清除),事后看来,都像是\dmark{}\origpage{1150}那场讨论会所发生之事的自然延伸,总
的说来也算得平淡无奇;那场讨论会的名字(或它的残余)也许会进历史,作为一次\textbf{警告}。
而翌年的“悼词”——无论在稍稍“摆一摆”它的人看来多么难以置信——同样像是这样一种延伸,
或者(如我先前所写\fnote{894}{(*)}{见注“宝石们”(n\textsuperscript{o}~170(iii))。})像一篇“尾声”。至于此后又过去的两年,它们不过是把一场光辉讨论会及其种种延伸
的“既得成果”,在文字里、在人心,一一盖棺论定……

这是一个显著的巧合——或者说,这分明\textbf{不是}“巧合”的效应——早在去年,在我还未见识
“SGA~4½--SGA~5”那场运算、也未见识反常层会议之前,我已经在我朋友皮埃尔与我的私人关系
中,辨认出恰好落在1977与1981这两个年头的两个“转折”。我第一次把它们纳入同一道目光、试
图探一探它们的意义,是在4月25日的注“两个转折”里——就在我发现\emph{葬礼}(读到那本难忘
的 LN~900)之后六天。两个转折发生的当时,早在多年之前,我远没有料到有一场\emph{葬礼}正在
暗中操办(至少在意识的层面上一无所知);要把其中任何一个同我所知的某件能照亮它的事联系
起来,我也会一筹莫展。

\subsection{殡仪馆——“为科学服务”}   % 18.5.8. Les pompes funèbres - “Im dienste der wissenschaft”

\notehead{175}{}

(3月23日)要把这场“葬礼运算”巡视完毕,还剩最后一位卖力的积极参与者的角色要看一看——
在这场关于葬礼的长考里,我已有机会屡屡“顺带”谈到他。这里说的是那家令人尊敬的Springer
出版公司(Springer~Verlag~GmbH,海德堡),科学书籍与期刊的出版者,大名鼎鼎——而且
还以“Im~Dienste~der~Wissenschaft”(“为科学服务”)的箴言而自豪\fnote{895}{(*)}{(6月1日)向Heinze博士打听之后得知,这其实算不上一句“箴言”,倒不如说是一条广告用语,英文作“Springer~for~Science”。}。这家出版社的数学书里,最负盛名的,大概要数“数学讲义
LECTURE NOTES IN MATHEMATICS”丛书。它恐怕也是世界上运气\dmark{}\origpage{1151}最惊人的科学丛书:二十来
年间出了一千多种。我想,我也为这份空前的成功出过一份力:这套丛书草创之初,正是靠我在六
十年代直至七十年代初,陆续放入的一大批我自己的或我的学生的著作给它撑的场子。我也曾以
“数学基础”(Grundlehren~der~Mathematik~und~ihrer~Grenzgebiete)丛书联合编辑的身份,与
Springer有过交往;三本书(包括 EGA~I 的重印本)经我的手在那里出版\fnote{896}{(*)}{另两本是Jean~Giraud与Monique~Hakim的学位论文(一是“场”与非交换 $1$-上同调的形式主义,一是任意赋环 topos 上的相对概形)。}。

1970年我告别数学舞台之后,我停止了一切编辑活动。仅凭惯性,我一直挂着那套丛书的编辑之
名,直到去年,才终于“正式”从Springer的一切编辑职务上退了下来。促使我这样做的是两个一
致的理由。其一,正当我重归“正统”的数学活动、重新发表数学著作之际,我要给这个“回归”划
出明确的界限:它对我绝不意味着重回某个“powerstructure”(权力与影响的架构),而只是回到
一份以发表为目的的\textbf{个人}数学工作。其二,自1976年起(经由Yves~Ladegaillerie
学位论文的那段插曲),我早有机会嗅到某种“葬礼的气味”——那时,对去年才发现的那场大规
模行动,我还毫无所觉。(关于这篇学位论文的插曲——它是我有幸启发过的最出色的论文之一——
见注“进步是挡不住的”(n\textsuperscript{o}~50),尤其是原委更详的注“二号棺材:或电锯式切块”,n\textsuperscript{o}~94.)这让我懂得:
“我所喜爱、我所愿意鼓励的那一类数学,在Springer出版公司里已经没有位置了”\fnote{897}{(*)}{这句引文(经翻译)出自去年2月18日致Peters博士的短信,信中我告知了他我退出“数学基础”编辑委员会的决定。Peters博士其实当时已经离开Springer(他如今在 Birkhäuser 出版社),通信便改由主管“数学基础”的Heinze博士接续。我曾要求把我那封信的副本寄给“数学基础”的每一位共同编辑(共十八位),并两次(1984年4月与1985年1月)向Heinze博士重申此请,他却始终认为无须告诉我副本究竟寄了没有(看来是\textbf{没有}寄)。我于是费神自己把信的副本一一寄给十八位编辑,并附了几句说明。他们之中我亲自结识七位,五位算得上我的朋友。只有一位(M.~Artin)费心回了信;其余诸位看来都没觉得有什么不对头——哪怕仅仅就\textbf{他自己}而言——:Springer公司居然没有费心把那封信(早在1984年2月)寄给他。};而且更重要的也许是,我在那里感到的气息\dmark{}\origpage{1152},并不鼓励我延续或重续与这家
公司哪怕稍微密切一点的往来。自去年2月我递交“数学基础”编辑委员会的辞呈以来,
过去的一年,只是把这份感受一再证实、一再加深。

不过这些都还是这场“葬礼运算”本体的边角——我昨天所说的那个“第二层次”,如今该回到它
了。据我所知,有\textbf{五本书}直接与这场运算相关\fnote{898}{(*)}{(6月1日)这几行写下之后,事情表明,还应在下面的清单上添入第六本书,它的书名本身就是一个骗局:“塔纳卡范畴”(Catégories~tannakiennes),Neantro~Saavedra~Rivano著。耐人寻味的是,这本书也出了Springer同一种“数学讲义”丛书。不过在那一桩运算里,Springer公司的责任在我看来并未牵连在内,不像其余五卷那样。至于“塔纳卡范畴”那场运算的细节,见那组注“第六根钉子(钉进棺材)”(n\textsuperscript{o}s~176$_{1}$--176$_{7}$)。}。按出版先后,它们是:SGA~7~I(1972年以我的名义出版)与 SGA~7~II(1973年
以Deligne-Katz的名义出版),整理的是1967/69年的单值群讨论班;题为“SGA~4½”(Deligne
编)的一卷,以及 SGA~5 的Illusie编订版(以我的名义出版,1977年);最后,是1982年以
Deligne--Milne--Ogus--Shih四人的共同署名问世、专为 Motif 掘墓复生立碑的那部“难
忘文集”。耐人寻味的是:这\textbf{五}卷皆出自\textbf{同}一家公司,皆在同一套“数学讲义”丛书\fnote{899}{(*)}{即第 288、340、569、589、900 卷。};前四卷出版时,主管“数学讲义”的是Peters博士\fnote{900}{(*)}{如我在倒数第二条脚注里指明的,Peters博士后来已离开Springer,转投 Birkhäuser 出版社。},最
后一卷出版时,主管丛书的是Byrne夫人。

这五次出版,其条件在我看来粗糙地不合规矩。我在别处已经指出:SGA~7~I 与 SGA~5 两卷
\textbf{以我的名义}于1972与1977年出版(LN~288与589),Springer却自始至终认为无须同我
打一声招呼,来征求我的同意,或仅仅把出版计划通知我。题为 SGA~7~II 与“SGA~4½”的两卷
的出版——它们打的也是 SGA 的旗号,而我以为这个旗号绝非人人可用,它与我的著作、我\dmark{}\origpage{1153}这
个人的干系,人所共知——事先同样没有征得我对使用这一旗号的同意,尽管(按理人们满可以期
待)我本该作为作者、或作为该卷(或它所整理的那个讨论班)的主编(或主编之一)列名其中
。最后,LN~900 一卷,在没有我的名分之下,呈现了那些概念、思想与构造;而凡消息
稍稍灵通的数学家都知道,这些是我引入的东西。在这一桩,明摆着(根本无须是 SGA~5 或 SGA~7
讨论班的那几个入室弟子之一)该卷就是人们通常所说的\textbf{剽窃}。我当然不指望Byrne夫人——出错
的话请指正——在该卷出版时主管“数学讲义”的那位——能凭自己的本事,看着稿子识破这场骗
局。但我想,一家严肃的出版社的本分,本该包括保证其出版物的严肃性,办法就是给自己配上胜
任的顾问。

这些顾问若诚实地做他们(想必)拿了报酬要做的事,也完全来得及向有关方面指出:SGA 这个
符号不是什么人都可以打的旗号,它自有\textbf{含义},而要尊重这含义,就该请教唯一有资格决定这个
旗号用法的那个人——也就是我自己。最后,就那卷顶着障眼法旗号“SGA~4½”的出版而言,作为
加重情节,只消翻一翻该卷的引论,或其后的那根“阿里阿德涅之线”,或第一章的引论,就能看
到:其中对待 SGA~4 与 SGA~5 两个讨论班,是何等漫不经心的轻蔑;而在消息稍稍灵通的人当中,
人尽皆知:那两个讨论班是六十年代中期的事,而顶着“SGA~4½”名目的一卷,装的却是七十年
代的伪作。依我看,对一个消息大体灵通、心智健全的人来说,这场骗局只能是一目了然的。因此,
更有十足的理由,在出版这样一卷、用上这样一个名头之前,先按正规手续征得我的同意。

依我看,在这五卷的每一卷的出版里,Springer出版公司的责任都牵连无疑——这五卷,正是围
绕我的著作、打着上同调主题所做的那场可歌可泣的诈骗大业中,一幕幕醒目的插曲。经由这些出
版,Springer公司成了这场非常行动的帮手与\textbf{贩运者}。我当然不能断言这是明知故犯。但我可以
说,这家公司在与我的关系上自\dmark{}\origpage{1154}1976年以来对我屡施的种种失礼(1970到1976年间,我想我没
有同它打过交道),也正朝着这场行动的方向,并且嵌入某种\textbf{精神},与它密不可分。

在注“操演”的子注“驱逐”(n\textsuperscript{o}~169$^{1}$)里,我曾提到我致Byrne夫人论 SGA~5 出版一事的信,和
她的复信——那复信,我必须说,真叫我目瞪口呆。(在这场“为科学服务”的光辉行动里,我“目
瞪口呆”既不是第一回,也不会是最后一回……)我从她的信里(日期是1985年2月15日)得
知:按照“一部著作包含多位作者的贡献时的惯常做法”(原文如此),无须专门知会我——我不
过是讨论班的\textbf{主编}罢了……SGA~5 的五位“作者”是Bucur、Houzel、伊
利耶、Jouanolou、Serre,独独不提我这卑不足道的人:我只以“主编”的身份列名——
想必纯属荣誉虚衔——为了这台光辉的讨论班,我说得未免太多\fnote{901}{(*)}{在那卷题为“SGA~4½”的书里那根著名的“阿里阿德涅之线”(贯穿 SGA~4 等等)之中,没有任何东西能让读者想到,我竟有幸在 SGA~4 与 SGA~5 里做过报告(相反,“有幸”为“SGA~4½”“合作”,倒是有的……)。关于这一点,见我在注“双重含义——或行骗的艺术”(n\textsuperscript{o}~169$_7$)里的观察,打字稿第899页。}。

早在收到这封富有教益的信之前,等得心里发闷(整整一个月杳无音信),我就拿出了我最漂亮
的笔(用德文),写信给 K.F.~Springer博士本人——公司主管负责人之一。那是一封漂亮的两页打
字信,我向他解释:公司与我的关系中绵延已久的一长串不快令我深为痛心;此外,还有若干起
对我不合规矩的粗暴做法,眼下我姑且奉上两桩,在我看来格外触目:“数学讲义”两卷(LN\textsuperscript{o}s~288、589)以我的名义出版,却认为无须征求我的意见;这两部书里,我在讨论班口头讲过的
思想、方法与结果,被删削,被残害,有时竟至于面目全非;这一事实,与环绕这两卷出版的种
种非常情形,凑在一起,对我而言绝不可能出于纯粹的偶然;我期待Springer公司公开地、毫无
保留地道歉,形式容当后议,一俟原则达成一致再行商定;我盼着他同我一样,把了结这一令人
不快、令人无法容忍的局面放在心上,找一个\dmark{}\origpage{1155}与事体相称的解决办法(“eine~dem~Fall~geziemende~Lösung~zu~finden”——这话说得还要更
雅些),落款则是规矩十足的“hochachtungsvoll”(谨上),出自我最漂亮的手笔。

开诚布公地说,我想我算是开诚布公了!这一回,Springer先生,总没法说,
没有人把情形亲自通报给他吧——而且是第一手的,通报者不是别人,正是当事的主要利害人
本人!

无巧不成书:回复(又过了一个来月)偏偏就在昨天到了我手上。它短得可怜,我禁不住要在此
全文照录(译文)。何况我还愣了一会儿,才明白这确实是上个月我那封漂亮信件的回复。回复如
下:

\begin{quote}
海德堡,1985年3月15日

亲爱的Grothendieck教授:

仍要感谢您2月9日的来信。Byrne博士夫人2月15日的来信想必已答复了您的问题。

此致,等等

K.F.~Springer
\end{quote}

这一回,至少我算彻底明白了!那班“消息灵通的人士”(前文已经说过的)想必跟他解释过:
何必为一个给他写信的、有点激动过头的先生费神——此人分明不属于上流社会。而且,倒也
不假……

还在等着收这封来自Springer殡仪股份有限公司管理层的、如此富有教益的回信的时候(承蒙
公司还肯屈尊,由大老板亲自署名作复,真是有礼),我有的是时间省察我自己的打算。这家可
敬的公司扮演的角色,在我看来实在太大;我一度设想过一场好看的大官司:以一位被冒犯的
“体面先生”的名义,作为无可名状的专横行径的受害者,索取天文数字的赔偿。可我又想:这样
一场官司,得吞掉多少精力。就算我赢了,并且拿到一笔叫人头晕的赔偿(往好处想!),那也是
若干年之后——我又落着什么好呢?我不缺钱,也不需要比我已有的更多;一场骗局,并不因为
某场官司打赢了或打输了,就多一分或少一分是骗局。我不会因此把世界变得更好,不会把我自己
变得更好,也不会把 K.F.~Springer先生和他麾下\dmark{}\origpage{1156}某些雇员的做派变好——当然更变不了他们对本
行的看法——即便我动员起我的律师,再让他们动员他们的\fnote{902}{(*)}{我还有过另一种念头:局面兴许会倒过来——这可敬的公司反过来告我一状,说我损害了它的名誉。这些“为科学服务”的人士,一碰到自己的名誉,想必是要较真的(只要涉及的是他们自己的名声……)。}。我也不会让某个上流社会里某种精神变得更好——那种精神,使得像Springer博士及其可敬公司(十三年来)甘为仆役的那类行动成为可能。我还剩下(我由衷希望)几年可
活——光阴似箭,而剩下的年月里,我满眼是叫人兴奋的事情。收集证据,好让法官们相信我与
SGA 有点干系——这类事想必不大叫人兴奋。我当年写它们,不是为了他们,也不是为了 K.F.~Springer先生,才费的力气……

至于那些(除我之外)我为之写作 SGA 的人:他们同这份东西——它(至少对我而言)仍是我
自身的一部分——所保持的关系,绝非与我无关。它是他们同我这个人的关系的一部分。说来奇怪,
这层关系我真正了解(或者说多少了解一点)的,只有我的五位上同调学生:正是托他们的福,
如今一位 K.F.~Springer博士才可以把我像块破布一样打发走,仿佛我对顶着 SGA 旗号的文字被如
何处置、或不如何处置,连发表意见的资格都没有——那位无名之辈的名字在不在封面上,都是
一样。

哪位数学读者若有耐心陪我读到此处,又若某一天也曾在 SGA(我说的是\textbf{真正的} SGA)里
出没,兴许愿意跟我说说他自己怎么看这件事。要是能收到一个字,来自某一位——他发觉,那部
我曾经独自一人整个儿扑在上面、花了十年生命、而\textbf{世上没有一个人}在工匠离场之后还肯费心
把它继续下去的著作——的确带着它的构思者与孕育者的印记,带着它成形之前那必需的全部岁
月的印记,并终于长成一间\textbf{人人可进的房子}\fnote{903}{(*)}{“盖房子”这个执念——而且是“对一切都好”的房子——在我的数学著作里扮演过可观的角色,早在五十年代初就是如此。这是我在劳动中的那股被我称作“服务的冲动”的具体表达;它属于那些赋予我的数学工作以生气的深层力量(尽管在“阴阳之钥”的思考之前,我自己都未曾察觉)。“房子”的原型第一次出现在我的思考里,是在11月26日的注“阴,作为仆人——与新主人”(n\textsuperscript{o}~135)里,来势之猛,出我意料。}——
收到这样一个字,我一定会高兴。况且,一间人人可进的房子,不是贫民区里一座公共小便池,谁
都觉得可以随意进去方便,在又破又黏的墙上涂写脏话……

\pmark\origpage{1157}而倘若读到这里的人,曾经是我的学生,或者曾经是我的朋友,却始终不觉得自己被
激发起写信的冲动、或当面谈一谈的冲动——哪怕只是为这件事,而不是为别的——那么让他知
道:他的沉默同样雄辩,而且,它是会被听见的。

\subsection{第六根钉子(钉进棺材)}   % 18.5.9. Le sixième clou (au cercueil)

\subsubsection{a.\ 预发掘}

\notehead{176$_{1}$}{}

\fnote{904}{(*)}{(6月16日)下面这一组注(n\textsuperscript{o}s~176$_{1}$至176$_{7}$),以“第六根钉子(钉进棺材)”为名,应当被看作“沉默”那组注(n\textsuperscript{o}s~168~(i)至(iv))的自然延续;那一组处理的是“Motif”运算,尤其是其中最后一道“预发掘”(n\textsuperscript{o}~168~(iv)),日期为4月8日。下面这些注——除最后一道(n\textsuperscript{o}~176$_{7}$)外——写于4月19日与20日。我之所以宁愿把它们撂在这里、“四种运算”的末尾,而不并入“Motif”运算,是因为此前几个星期里围绕其余三种运算——尤其是那场(所谓“反常层会议”或“当值无名之辈”的)运算——继续进行的思考,而那场运算正是“封神”那组注的题目——给刚刚出现的那个(同样出人意料的)“新事实”,投下了意想不到的光。这里提醒一句:写下面这些注的时候,原则上,我已经在\emph{葬礼}底下画过“句号”(其终极之注“朋友”(n\textsuperscript{o}~188)日期是4月7日),并且打算一两天内就把葬礼(III)的全部手稿送去打字。也就是说,这些笔记是在“最后一刻补遗”的心境里写下的……}

(4月19日)我终于有机会(4月10日)读到了“预发掘”一注(n\textsuperscript{o}~168$^{1}$)所引的Langlands的文章。照Deligne去年8月寄给我的那份 Motif“带评注书目”的说法,Langlands这篇文章,与Deligne刊在同一卷里的那篇(即前引之注的对象)一道,是自1970年我离去以来头一批用到 Motif 的文章\fnote{905}{(*)}{不过,1972年Kleiman与Saavedra的那些报告除外——它们承接了围绕 Motif 范畴描述的那几份微薄的“练手音阶”(可参照注“把一个个 i 点清”(n\textsuperscript{o}~164)里打字稿第794页脚注(**))。}。直到去年我还不知道Langlands的文章(Deligne的那篇也不知道),这情有可原:作者(同我的前学生一样)不认为有必要给我寄一份抽印本。人们也不禁要问:他何苦费这个神呢——翻一翻他的文章就清楚,我这卑不足道的人,同他文中所谓“自守表示、Shimura簇与 Motif”那个题目,严格说来毫无关系\dmark{}\origpage{1158}。我的名字(借用一句我的打字机恰好一年前就背熟的套语)在这篇文章里、在书目里,一处也不出现。不过,有几个我大约在1964年提炼的观念(或者是我梦见自己提炼过——看来我还在重蹈覆辙……),我自以为认出来了;我甚至把这个梦的追忆(或者也许是一个关于梦的追忆的梦……)白纸黑字写了出来,恰恰就在1984年4月19日\fnote{906}{(*)}{关于这一点,见注“一个梦的回忆——或Motif的诞生……”,n\textsuperscript{o}~51。}。我恍然觉得自己回到了整整一年前的同一天。

诚然,隔了这一年,我已有时间见怪不惊。纵有不快,也几乎算不上意外(人家会说,本来也没多
少东西……),更谈不上震动。何况,那本为三年后那部难忘文集 LN~900 鸣锣开道的先导文章,与
后者之间有一个不小的差别:我没有荣幸亲遇Langlands;他的“Motif Galois群”(或“基本群”)的瑜伽,
不是从我嘴里听来的(Deligne1965或1966年前后可是听过的)。不过,整个六十年代下半叶,我在
周围、在凡愿听者那里,把它讲了个够(Langlands总不是刚下船的外乡人……),所以我有一份揣测:
Langlands心里明镜似的,知道这门关于各种Galois群与基本群的新“几何”哲学——把它们看作合宜的
仿射 pro-代数群——是从哪儿来的。我猜他心知肚明:这门哲学,并非1972年从某位Neantro~Saavedra~Rivano的脑子里生出来的——此人此后从流通中销声匿迹,没有留下任何痕迹\fnote{907}{(*)}{据Deligne去年10月来访时向我暗示的,Saavedra实际上已经改了行(如今“搞经济”了),自1972年论文答辩之后,就再没碰过数学。}。我觉得,Langlands倘若觉得有必要,就此表个态,并不算什么奢侈。当然,照眼下的世
道,盼他费这份神,也许是我过分乐观了……

\subsubsection{b.\ 意外的惊喜}

\notehead{176$_{2}$}{}

\pmark\origpage{1159}好事从不单行。读到Langlands那篇文章的次日,我又有机会翻了一翻Neantro~Saavedra~Rivano的那本Langlands大量援引的书,书名叫“塔纳卡范畴”(《数学讲义》265卷,1972)
。

我离去之前有过九个(男)学生,Saavedra是其中唯一一个我后来再无音讯的;也因此,没有任何
风声向我提示,他多少沾上了某场\emph{葬礼}的“颜色”或“气味”。我于是草草下了结论——带着我惯
有的那份天真:他(哪怕只因为没了机会——听说他后来离开了数学的水域……)是所有学生中唯一
全然置身于“葬礼运算”精神之外的人。可是,同Jouanolou的情形一样,音讯稀少到这个地步,本该给
我提个醒。我当然知道:他跟我共事时预定要成为论文的那份东西,最终以“塔纳卡范畴”之名,于
1972年在“数学讲义”出了书;而那本书,上星期之前,我不记得自己曾费神看过一眼\fnote{908}{(*)}{(6月16日)Saavedra想必不认为有必要寄我这本书;我没有它的本子。不过,七十年代里,我也许亲手捧过它。我留着的印象,仅此而已:他做了一份细致周到、拿来就能用的活儿;至于这印象从何而来,我说不准。这印象尤其在场,是在写“清场”那注(n\textsuperscript{o}~67,尤其打字稿第252--253页)的时候;那里我评说了Deligne的这个“谜”:他几乎逐字“誊抄”了Saavedra跟我一起做的那篇论文。}。我满脑子别的任务,压根没想到:Saavedra从此杳无音信,连他的论文答辩也不
通知我一声、不请我进答辩委员会——按说我是最清楚底细的人——这有点奇怪。正是读了这本书,
才明白他为什么宁可不打扰我,而把论文“偷偷”过了关,交给一个我完全不知其组成的委员会
\fnote{909}{(*)}{这个答辩委员会组成的谜,在“第六根钉子”的第七道即最后一道注(n\textsuperscript{o}~176$_{7}$)里,以全然出人意料的方式解开了;此处不再多置一词……}。那时\emph{葬礼}已经在加速进行:因为委员会里没有一个人认为有必要通
知我答辩一事,更不用说请我参加委员会了(Jouanolou的论文答辩,大约也在同期,情形如出一辙)
\fnote{910}{(*)}{关于一次更正,见前一条脚注所引之注。}。

\pmark\origpage{1160}这本书,呈现的是那门“算术几何”的一个关键侧面——它的视野(总体眼光)六十年代里在我心中诞生
并成长(当时还没有名字),而 Motif 的瑜伽曾是(并且始终是\fnote{911}{(*)}{不过,其间,这份“灵魂”又添了“非阿贝尔瑜伽”的滋养;《纲领概要》里略微谈到过它。(关于这份文本,见《引论》之三“罗盘与行囊”。它反正将收入《数学反思》第3卷。)})它的灵魂。可以说,Saavedra这本书的实质,是对我的某些
思想的一份形状完整、细致周到的陈述——关于某些范畴的一门“Galois--Poincaré理论”(我做梦也
不会想到叫它们“塔纳卡的”……);那些思想,是我在Saavedra还拿不准肯不肯下功夫去熟悉、去
消化,以便能把它们收进其论文的“陈述性”部分的时候,长久的、不厌其烦地讲给他听的。我曾交给他
一批原原本本的手写笔记,命题钉是钉铆是铆,证明概略俱全;至今我还等着他肯把它们还给我
\fnote{912}{(*)}{把我的手写笔记这里那里散发给学生,按需取用,是我的老习惯——而他们要学的头一
课,就是辨认我的笔迹。向来讲好:他们用完,就把笔记还我——但据我所知,这个愿望很少被遵
守。这不过是众多征兆之一,表明我的学生并不怕我,而宁可把我看作一块“好脾气”的面团:干起活
来要求是严的,除此之外,随和得没边……}。当然,论文本身的题目,并不是去陈述一个其动机他完
全无从把握的人的思想。要做的,是显式地给出我在这里称为“Galois--Poincaré”范畴的那些“张量”范畴\fnote{913}{(*)}{为了不叫它们“Grothendieck范畴”!然而,在我有幸引入并命名(因而并不冠我之名)的众多范畴与新观念之中,若有一类非用这称呼不可——单凭体面,我就想写——正是这一类!(topos 除外,但它的名字我觉得本就好……)至于“塔纳卡范畴”这个名字——由一位天才的前学生悄悄塞进来的,又被全体一致的教团欣然采纳——这不多不少,正是一场骗局,如下文原原本本所述。(见“会等待的人……”之后的注,n\textsuperscript{o}~176$_{3}$。)}的一种“便利”的内蕴刻画,亦即:允许对它们作“Galois--Poincaré--Grothendieck式”的描述,把它们表达为基底环 $k=\mathrm{End}(1)$ 上一个“仿射(pro)代数 gerbe”的线性表示。当 $k$ 是域时,我曾以(我引入的术语)所谓“刚性”性质给出过这样一个条件;我记得自己在1964/65年、思考 Motif Galois群之始,就写出过一个完整的证明\dmark{}\origpage{1161}\fnote{914}{(*)}{我没有费时去我的笔记里核实这件事——关于 Motif Galois群的笔记(或者说,其中没交给Saavedra的残存部分)。反正第3卷里我要回去谈,多半在“Motif,吾爱”那一章里。}。我得把原理指给他;至于我在这题目上写下的笔记,我克制着没有交给他——因为
学他未来的行当,是他自己的事,不是我的事,活儿得他自己做。如果我记得不错,当时对我而言悬
而未决的,只有一件事:这样的Galois--Poincaré式理论,就基底环 $k$ 须满足的假设而论,其自然有
效域何在——我尤其关心 $k$ 取 $\mathbb{Z}$ 这样的环的情形(因为牵涉到 Motif 理论的应用)。

在我离去前的所有学生里,Saavedra是来得最晚的\fnote{915}{(*)}{如果我记得不错,Saavedra是1968或1969年要求跟我工作的,在我不期然告别数学舞台之前一两年。},也是准备最差的,而且(至少起初)最没
有干劲想“露一手”。所以我不抱希望:他会超出我给他出的那个非常有限的技术问题——它只要求最
微薄的知识(一点概形语言、线性代数、平坦下降、gerbe 的语言,如此而已)。而他书里第四
至第六章所处理的那些更精细的问题(纤维函子的滤套、实域上Galois--Poincaré范畴上的极化结构,
以及“可极化”的这类范畴的清单,及其对 Motif 范畴和许多变体的应用),要求的却是“四面八方”
的知识\fnote{916}{(*)}{这尤其要求对约化代数群结构理论、及其在实域上的分类的深入了解,还要熟悉整整一批观念:Motif、晶体、$F$-晶体、分层模、局部系统(而此人当初连拓扑空间的奇异基本群也只略知皮毛),再加上 Hodge 理论,以及那些从未在文献里显式陈述、只“藏在行间”的精微的“极化”性质。},从而要求可观的入门功夫;我不认为Saavedra供得起。我至多盼着他也许在论文后
面附一个摘要(多多少少由我口授),把那套理论中不会收进形状完整的陈述性工作里的要点列一
列。直到上星期我才被纠正过来,并且意识到:Saavedra交出了一份真正惊人、而且是以破纪录
的速度完成的活儿\dmark{}\origpage{1162}\fnote{917}{(*)}{关于这个“纪录”的更深的思考,以及它(明摆着的)解释,见注“Monsieur~Verdoux——或献殷勤的骑士”(n\textsuperscript{o}~176$_{5}$)。}。这活儿结晶为一本书:陈述细致周到、一丝不苟,分毫不差地
摹拟着、拿来就能用,以(在我看来)近乎穷尽的方式,呈现了我在六十年代发展的那套几何-代数
形式主义。从这个角度看,我认为他做了一份有用的、处处都无愧于人的工作;我方才说的那份“惊
讶”,确实是一份“意外的惊喜”。

这份工作,精确地说,在于把我提供的一整套观念、命题与证明,以“典范的”、可径直发表的形状
整理出来(按照我当年那些严格的标准)。这样的陈述性工作,当然属于数学家这一行的本分——无
论处理的是自己的观念与结果,还是他人的。与我许多同事不同,我不认为:在评估一篇论文或任
何出版的质量时——及至授予做这种工作的人数学“博士”的头衔、把他当作不折不扣的数学家
时——这样的工作可以被当作可忽略的量。不过,有一点在我看来至关重要:这一行的某种起码的
伦理必须得到尊重;凡工作在于陈述与发展他人观念之处,这一点必须清清楚楚地写明,不留丝毫
暧昧。

然而,在眼下这个案子里,全书除了一篇光辉引言末尾附着的、三行纯客套的、含混其辞的“致
谢”\fnote{918}{(*)}{这篇引言的实质,是逐字照抄四条主要命题——那是我向Saavedra指出的、有待发展的Galois--Poincaré瑜伽的“四根台柱”(至于纤维函子上的滤套所牵涉的问题,它们很难被概括成一条石碑式的命题,故未列入);但其中一条——本该构成其论文“最低纲领”的那条——被添加了一个显赫而明摆着的错误,使得它平凡地成为假的!下一道注(“会等待的人……”,n\textsuperscript{o}~176$_{3}$)就会谈到,尤其见已引的注“Monsieur~Verdoux——或献殷勤的骑士”(n\textsuperscript{o}~176$_{5}$)及其后一道注“肮脏的差事”(n\textsuperscript{o}~176$_{6}$)。}之外,没有任何东西能让读者猜到:书中展开的任何一个主题——从全书的主题本
身开始——跟我这卑不足道的人有什么相干。我恍然又回到读到那本纪念性的 Motif 掘墓文集的
那一天(到今天恰好整一年,一天不差)!我的名字在全书里几乎无处可见,只有两三处例外:那里
需要形状完整的参考文献,而除我之笔,别无可用者\dmark{}\origpage{1163}。

而且,这绝不仅是——怕被人看清作者“只是”在陈述别人的观念与结果而感到的那点尴尬(这——
尤其在这个案子里——若工作做得聪明,已经很不少了)。我从许多不会骗人的“小细节”里看清:这
里绝不止是稍微“收割”一把、给自己的纹章镀点金、然后退回幕后的意思。这真是\textbf{为葬礼而举
行的葬礼}。只举\textbf{一}个例子——上帝知道,我曾花了多少天、多少个星期,向Saavedra原原本本、
翻来覆去地讲解那些观念:他初来乍到,一无所知;晶体、$F$-晶体(在特征 $p>0$ 的情形顶替缺失
的 $p$-进“系数”,好定义 $L$-函数……)、分层模(及其与局部系统的关系),最后还有一点 Motif
的瑜伽(暂时以标准猜想作启发的底座);这一切,为的是借一大把例子让他明白,我拿这些Galois
--Poincaré范畴要到哪里去;也为的是——以防(谁说得准呢……)他有勇气有毅力,在预定的“最低纲
领”之外,至少收进一章典型的例子。他心里一清二楚,不必我多费唇舌:这些都是关键的几何观念,
而且并非开天辟地就有;把它们讲给他听、一遍遍讲给他听而不知疲倦的,不是别人,正是我——
它们是我在此前五到十年间引入的,用来伺候某一视野(总体眼光)的工具(即使那视野从他头顶掠过,
正如它掠过我所有学生的头顶,只除了一位\fnote{919}{(*)}{此人在师父一转身之后,就按捺不住把它埋葬了……})。可是,在他引入并稍稍发展这些观念的地方(在专述例子的第六章里),也如同在他佯装以作者自居的那部分理论正文里,我的名字照样不见踪
影。然而我很难设想,Saavedra会真的以为:读者无论多么孤陋寡闻,哪怕满心愿意相信他是这些
范畴(他慷慨地称之为“塔纳卡范畴”)的父亲,也不至于进一步以为:连 $F$-晶体、Motif 与“塔纳卡
装备”(原文如此)里的其他小玩意,也是这位Saavedra为应一时之需发明出来的。倘若这些观念被
当作信手拈来的寻常货色、或从最近的\dmark{}\origpage{1164}孤儿院里捡来的弃儿,那我在其中认出了一种\textbf{风格}——
一年来我为\emph{葬礼}兜圈子兜个没完,这种风格我再熟悉不过了……

是Mebkhout把那本书带给我的,他满心欢喜,总算能给我看一个我的学生的例子:这位至少还是
“诚实的”\fnote{920}{(*)}{(6月16日)再说,这事没成,他可是难过极了,还尽力来安抚我——这让我想起Kawai的情形(见打字稿第1078页脚注(*)),也想起了Beilinson:Mebkhout觉得Beilinson“比Bernstein更诚实”(见第1072页)——真像第欧根尼提着灯笼,只不过这一回,是在那伙一头扎进上同调这个名声实在不佳的题目的人当中,找一位诚实的数学家……}。他显然被引言末尾那三行致谢打动了——也难怪:1972年,感
谢某位已故者的风气,本来就已不大流行;从那以后,在我不止一位前学生那里,取而代之的调门,
不是冷嘲热讽,就是打趣,再不然就是彻底的沉默。话虽如此,这一回我到底得到一份“深深的感
激”:感谢我“把[作者]引入了这个题目”,感谢我的“劝告与鼓励……对完成这项工作必不可少”
\fnote{921}{(*)}{这些“致谢”,就情形而论,不过是一句玩笑话:人们会以为,我把作者“引入”了单复变函数这一“题目”,或任何同类的经典题目。其实,当我在一位正愁没论文的Saavedra面前谈起这桩事时,这个“题目”还并不存在——除非是在我之内、与 Motif 共生而发展出的那一份视野里,以及在我赋予它形状的手写笔记里。关于这份视野的诞生与发展,我在“一个梦的回忆——或Motif的诞生……”一注(n\textsuperscript{o}~51)里谈到;至于我的一位前学生(在众人纵容的目光下)对这根源所作的、漫不经心的彻底抹除,见其后一注“葬礼——或新爸爸”(n\textsuperscript{o}s~51、52)。(6月16日)Saavedra的这些致谢之所以更是一句“玩笑话”:其作者从来不曾费心,哪怕只寄我一本他的书、连同这些冒牌的致谢。既然我已经把“塔纳卡范畴”(原文如此)这场运算巡视完毕,如今我更明白:我这位前学生没有理由为他的“工作”(原文如此)感到骄傲,也不急于让我看到它。而照当时、乃至直到两年前事情的走势看,那位工匠似乎根本不大可能有机会看到它……}。这就是所谓的拿话头搪塞人;而在我看来,在\emph{葬礼}分明大行其道的时刻,在陈述
自己的工作时保有起码的诚实,才是表达“感激”的更有力的方式。
\subsubsection{c.\ 会等待的人}

\notehead{176$_{3}$}{}

\dmark{}\origpage{1165}说实在的,这本书一到手,我就看清:在那场难忘的“SGA~4½--SGA~5 运算”之前,整个
\emph{葬礼}里,还没有哪一幕的分量,赶得上这本名叫“塔纳卡范畴”——一个再平淡不过的书名——的
LN~265。此前的插曲\fnote{922}{(*)}{这里说的“插曲”,在“埋葬……”那条注(n\textsuperscript{o}~168~(ii))里匆匆勾勒过;它属于“Motif”运算那组注。}不过是或隐或显的“收割”,把某些重要思想的承续隐瞒起来。而这
里,我的“算术几何”视野的一个整个的关键侧面,被人不动声色地“挪用”了——经手的,竟是我的
所有学生当中看上去最“不起眼”的一位!

\pmark 的确,在这位学生背后,我凭一种不会骗人的笔法,认出了那个牵线的人;而且此公在Saavedra极尽殷勤的致谢名单里,赫然排在前列\fnote{923}{(*)}{就“数学”本身而论,这几位依次是(按出场顺序)我自己(不按字母序,算他有礼貌)、Berthelot与Deligne。}。Saavedra给这部书、给我引入的那个关键观念所起的名字,是一件精巧
的\textbf{褫夺之作}。论其斩钉截铁的效力,五年之后才被另一个名字超过——这回,那名字出自Deligne
本人之笔,落在另一部书上\fnote{924}{(*)}{正如下文(已引的注“Monsieur~Verdoux——或献殷勤的骑士”)将显现的,至少有强有力的揣测:与其读作“但这回出自Deligne本人之笔”,不如读作“并且同样出自Deligne本人之笔”。}。

\pmark 如果说,顶着一卷“锯倒之作”的“SGA~4½”之名,是一桩天才的冒名顶替,那么“塔纳卡范畴”
这个名字,就是一桩同样天才的\textbf{骗局}。哪怕对一个“平凡”或“中性”的Galois--Poincaré范畴——
等价于域 $k$ 上一个仿射群概形 $G$ 的有限维线性表示的范畴——我发展出的那套瑜伽,也十足是
“Grothendieck式”的:它由我早先在拓扑空间、概形或(更一般的)topos 的基本群情形发展的类似
瑜伽启发而来。把基本群定义为“覆盖层范畴上一个纤维函子”的自同构群,以及(同样古怪,因为
新,故而不惯)系统地同“不必\dmark{}\origpage{1166}\textbf{连通}”的平展覆盖的范畴打交道的想法,当年没少招来讥讽。我从
来不放在心上,深知那班爱开玩笑的家伙——自以为在学堂板凳上学过Galois理论或Poincaré理论就
算懂了——没有一个真正懂过;而且直到今天,他们之中谁也不能,不逐字照抄我做过的东西、不照
搬我以范畴等价给出的覆盖的Galois--Poincaré理论的表述,独自迈出(比方说)稍一般些的概形的覆
盖的Galois理论的\textbf{最初几步}\fnote{925}{(*)}{“稍一般些”在此可以精确理解为“一个非正规的概形”。在我之前,代数簇的基本群(由朗与Serre)只在正规簇的情形引入过,办法是把它描述为其函数域的“绝对”投射有限Galois群 $\mathrm{Gal}(\overline{K}/K)$ 的一个合宜的商。}\fnote{926}{(*)}{今天,这种表述基本群与覆盖之关系的方式——哪怕在普通拓扑空间(局部弧单连通)这个可以说是“教科书式”的特例里——也开始到处流通,而对我们共同的先祖只字不提,这还用说么……}。

同样,从一个“抽象”的有限维线性表示范畴出发——带着它自然的乘法结构与它自然的“忘却 $G$
之作用”的纤维函子——重建一个(为确定起见,取域上的)仿射群概形,把它作为\textbf{这个函子的自
同构的群概形}——这个想法,既不归功于塔纳卡(他可从未要求这么多),也不归我那微末的前
学生Saavedra,也不归我最出色的学生Deligne(令我遗憾——不过那时他还没在跟前),而是典型
“Grothendieck式”的。由此得到 $k$ 上仿射群概形与配有 $k$ 值纤维函子的刚 $k$-张量范畴之间的完
美对应,道理亦然。再者,倘若碰巧(特征非零域上的 Motif 范畴,就往往如此)一个刚张量范畴
(不幸,或者幸运加倍……)偏偏\textbf{没有}纤维函子可用,那么“代数群”就该换成一个代数
“gerbe”——这个想法,早在年轻的Deligne在数学里还没听人说过“gerbe”这个词、做梦也不曾想
到会有类似东西的年代,就已经被阐发得淋漓尽致。同样,当Giraud在六十年代自告奋勇,发展一整套
二维以内的非交换上同调代数的武器库\dmark{}\origpage{1167}——栈、gerbe 与 lien 的语言\fnote{927}{(*)}{这套形象化的术语是Giraud引入的,代替我从1955年起使用的临时术语(有点随手抓来,如“局部性质的纤维范畴”之类不成体统的名字;那些观念的根本性质,要求的是斩钉截铁、掷地有声的名字)。
(6月16日)在他那本书的引言第一页,Saavedra说到了Giraud“引入”的非交换同调代数“形式主义”。这是众多地方中的一处:在其中,我感觉到有一个比本书作者更机灵的人,在“搀着他的手”……正是同一位,一提“导出范畴”,就乐于紧接着补一句“由Verdier引入”(而这两处的底细,他心里一清二楚……)。}——讥笑声同样不缺。这类东西,如今、而且早就被Deligne之流叫作
“废话矿渣”。这些讥笑不曾让我却步\fnote{928}{(*)}{Giraud却不然——他在刚与我一同开辟这主题不久,就与之决裂,再无回头。关于这一点,见“共同继承人……”一注(尤其打字稿第386--387页)及其后一道注“……以及那把电锯”(n\textsuperscript{o}s~91、92)。};我知道自己在往哪里走——我看着这块“矿渣”以
完美精细的手法,抓住那些微妙而深刻的关系(我深知,任何别的“语言”都无力把握它们),心中是
“复苏”(如我在别处所写),却谈不上惊讶。

话虽如此,当这班讥笑者某一天发现一道被自己漏掉的“奶油甜点”——无论是那被人抢着命名为
“塔纳卡的”(且听下回分解……)范畴,还是某种“对应”或“关系”或“构造”(边上带点新Grothendieck味儿),被人用委婉语一笔带过、或被命名为“Riemann--Hilbert”(同样且听下回分解……)
\fnote{929}{(*)}{关于这后一个“且听下回分解”,见“反常层会议”的全套,尤其注“魔术师”与“冤大头买卖——或木偶戏”(n\textsuperscript{o}s~75$''$与 171$_{2}$~(e),后者属于长注“黑手党”(n\textsuperscript{o}~171$_{2}$)的一部分)。}——这时人人蜂拥而至,争当光辉的发明者。这就是本世纪七十年代、八十年代数学的
“时代精神”……可以确定的是:给这些范畴(我曾长久讲给他听的)起“塔纳卡范畴”这个真正光辉
的名字,这样的主意,不是一个Saavedra想得出来的。放任他自己,他绝不敢改动从我这里承袭的
术语,至少得先征求我的同意——这原是最起码的规矩!非得有上头的榜样与鼓励,他才敢\dmark{}\origpage{1168}这样把我
当作可忽略的量打发掉。再说,这可怜人为了熟悉必不可少的东西——若他想完成我交给他的那个雄心勃
勃的写作计划的哪怕一部分——手头的活儿已经够多了\fnote{930}{(*)}{他以仅仅两年的创纪录时间完成了这个计划——从我离去(那时计划几乎尚未动工,除了对概形基本技术的初步熟悉)算起。即便有Deligne在一旁搀扶(在我离去之前,他对这个学生毫无兴趣),这份成绩也简直堪称奇迹;这个“奇迹”,在注“Monsieur~Verdoux——或献殷勤的骑士”(n\textsuperscript{o}~176$_{5}$)里得到较细的检视。},哪有余暇再去翻文献、去读塔纳卡以及天知道什么的东西——他跟我共事的时候,肯定
从没听说过这些\fnote{931}{(*)}{提醒一句:Saavedra跟我工作,恰好在我离去前一两年(约1968、1969年),此后我几乎完全失去了他的踪迹。他那时的行囊,同任何一个来自第三世界(或我们外省院校)的三年级研究生相比,不厚也不薄。}。

这个名字之所以“光辉”,在于两种看似矛盾的品质的精巧结合。其一,在肤浅的观察者眼里,这
名字不显得全然荒唐。“人人”都约略记得存在一个“塔纳卡对偶”,乘法结构在其中扮演角色——
这看起来同那著名的 $\otimes$-范畴(某位Saavedra——这是谁呀?——称之为“塔纳卡的”)身上发生
的事,还真有几分像;那就“塔纳卡的”吧,为什么不呢!可是,对会等待的人,事情自己会成熟。十三
年过去;如今,顶着谁也没见过的无名作者之书的,换了一部三年前问世的、显赫得多的参考:光辉
的 LN~900,出自非别人而是Deligne之笔,外加一位名叫Milne的搭档。这两位大名鼎鼎的作者,从
头发展了他们同样称作“塔纳卡的”那类范畴的全部形式主义。显然,这是一个根本性的
观念,多年来已被Langlands、Deligne、Serre等人使用,前程似锦。当然,谁也不会相信:在这篇文章里被
顺带引了两三次的某位Saavedra,竟是这一关键观念、以及由它生出的那套精微形式主义的主人。
两位光辉作者的这篇文章,以其主笔人所熟知的全部娴熟,重拾这个题目,其笔调本身,就把这层意\dmark{}\origpage{1169}
思交代得没有半点含糊\fnote{932}{(*)}{关于这篇文章,尤其见“下葬——或新父亲”(n\textsuperscript{o}~52,尤其打字稿第214页)与“清场”(n\textsuperscript{o}~67,尤其打字稿第252--253页)两注。}。更不必说,他们还从Saavedra书里陈述的理论中,挑出一处粗糙已极的错误(以致
他们不得不从一个全然不同的定义出发——那个定义,终于像是正确的\fnote{933}{(*)}{关于Deligne的这桩绝技(Milne作陪衬),见屡次所引注“Monsieur~Verdoux——或献殷勤的骑士”(打字稿第1176页)的开头。}),叫人不免要问:这可怜的
Saavedra(当年必有某人——猜猜是谁……——费过心给他解释是怎么回事)到底懂不懂自己在
说什么。而Milne纵然光辉,且有幸与显赫的Deligne合署一篇文章、发展一个分明根本性的观念,
也决不会起念头认为自己能算它的父亲或哪怕共同父亲;Beilinson与Bernstein同样不会声称,他们发
明了(或哪怕共同发明了……)那句有名的“本该在这几篇笔记里找到自己位置的关系……”——他们
有幸与同一位显赫的Deligne合署了它,而在那之前,正是Deligne以他无上的好意,给他们指点了
Kazhdan--Lusztig 证明的门径……最后,又有谁会\textbf{认真}相信:那位把自己的名字(想必未经商
量)借给这一根本观念的、大名鼎鼎的塔纳卡,果真与此\textbf{有}几分干系?等到哪天世人都看清了这
个观念、以及随之而来的整套精微理论的\textbf{真正的父亲},他也不至于出面认领——倘若他还在世
的话。对此若还有半点疑虑,只消翻翻塔纳卡的著作——嫌多的话,就翻他论“塔纳卡对偶”的那
篇——就会发现:说到底,风马牛不相及……

这里,仍然是:几块界石一立,剩下交给时间。显然,这门理论——它将越来越多地显形为一种新的
\textbf{哲学}的技术手段,用以沟通几何与算术——在未来的年月里,注定要越来越走到数学舞台的中
央。五年、十年之内,再不会有人想到,就这个题目去援引某个无名作者的书;而那位想必搀过他手
的人,倒是不辞辛劳,写出了一份势在必行的陈述,由一位光辉的合作者协助,构成那部光辉的文集
的核心——Motif 的观念,终于在\dmark{}\origpage{1170}那片坚实的土地上得到发展。(至于那位含混而潦
草的先行者——他早已被遗忘——以及那道分明超出他能力的、惯常的猜测性“废话矿渣”,文集里
只字不提,倒也显得仁慈……)引用“P.~Deligne 与 J.~S.~Milne 的塔纳卡范畴”,将如同引用
FAC 或 GAGA(Serre的)或 SGA(法国高等科学研究所那部著名的匿名讨论班,人称“布瓦马
里”)一样,成为第二天性。而这么做的时候,任何人的心里,对这些新观念的父亲是谁,都不会有
半点含糊——那父亲断然不是合著者Milne,更不是塔纳卡,甚至也不是某位严格无名的作者(一个
叫Saavedra的):他在人家文章里被顺带提了两三次,为的是他(在他那卷书引言里)写过关于这
个题目的一份“出色的摘要”(带若干保留)。

不过,人们不会指望这门理论的父亲,让他去强暴自己那份出了名的谦虚,以至于把(按“消息灵
通”人士——他们在这类事情上说了算——的全体共识,理应如此称呼)本该如此称呼的东西,
叫作“Deligne范畴”(或者,在另一个天地里,“Deligne对应”)……

\subsubsection{d.\ 父亲们的圆舞曲}

\notehead{176$_{4}$}{}

(4月20日)昨天的思考,让我以新的眼睛看清一件去年——我刚一头栽进\emph{葬礼}的时候——曾让
我瞠目的事:“……这件看似荒唐的事:Deligne十年之后‘重做’Saavedra的论文*!”从去年4月19
日——我读到那本“难忘文集”LN~900 的日子——起,这件事就被谈到了:文集里(除其他美好东西之
外),Saavedra的论文被几乎逐字重现\fnote{934}{(*)}{见倒数第二条脚注所引诸注。}。一周之后,我在“清场”那条注里又回到它。当时我已得出“内
心确信”:这堆废话背后的意思,是光辉的Deligne(自任Saavedra的文书)的如下愿望:
\begin{quote}
“给自己一种幻觉中的解放感,解脱他肯定感到的一桩痛苦的义务:不得不时时援引那个他正要
取而代之、加以否定的本人,或哪怕任何一个援引了他的人。”
\end{quote}
可上星期,我第一次费心翻阅了这位“别人”的工作,惊讶地发现:他压根\dmark{}\origpage{1171}没想“援引我”(除那三行
冒牌的“深深感激”之外——明摆着是障眼的烟幕)。于是,一年前我的“内心确信”就站
不住脚了——里面肯定有对的成分,可仍剩下谜:总不至于就是那三行字(没有哪个读者会想到跑去
引言末尾把它们挖出来),竟能促使一个Deligne,去当一位早已作古的宗师门下最无名学生的抄写员
吧!更不必说,在这段引言的末尾,我几乎与他、与Berthelot并列出现:他们同我一样(人们会说
\fnote{935}{(*)}{只是有这样一点不同:是我“把他引入了这个题目”(原文如此),而且“他至今在数学家的养成上,还大大受惠于我”(这真是受之有愧)。}),有资格得到“他们在这一工作中慷慨给予的帮助与建议”的
致谢……

这个“谜”在昨天的思考里完全解开了;我不必去找,甚至不必刻意去想。写完之后重新回味,种种联
想浮现出来——写作之时它们必定已经在场,只是我毫无察觉,还在暗中引导我的笔。打动我的,不
仅是风格的相似,更是一整套\textbf{专利的据为己有之术},贯穿\emph{葬礼}的三大“运算”(四者之
中,Deligne本人是主要的——即便不是唯一的——“受益人”)。不妨称之为“临时替身爸爸”的手法:
悄悄把它摆上数学勒索的棋盘,好变走一份真实的父亲身份,而我朋友皮埃尔其人,则暂时隐在幕
后。一旦生身父亲被体体面面地彻底清出舞台,替身爸爸也就被变走,仿佛从来不曾存在;于是\textbf{真
正的父亲},谦和地微笑着登台,甚至无须亲口说破——因为那个不动声色牵线、又懂得等待的人,
诸事自会水到渠成,毫无阻拦:教团全体的同心一致,早已把他按他所应得的名分加了冕。

这套手法,直到前几天,我循着我朋友Zoghman在第四场运算(所谓“当值无名之辈”)各插曲中的
倒霉遭遇一路追索,才开始看清。这一场里的“替身爸爸”(为某件“对应”……)是\textbf{Kashiwara}——我
说不好,他是就这么从天上掉下来的,\dmark{}\origpage{1172}似有天助且纯属巧合,还是未来的真爸爸巧妙地让他明白:
一个无名之辈的那个没有像样父亲的结果,横在这里,弃之可惜\fnote{936}{(*)}{(6月16日)看来,对Mebkhout著作施行扒手手术的主意,正是企业家Kashiwara出的,而且早在1978年——其时Mebkhout刚把刚杀青的论文第三章拿给他看不过几个月。关于这一点,见注“黑手党”部分(b)(“最初的麻烦——或太平洋彼岸的黑帮老大”),打字稿第1060页脚注(*)。}\fnote{937}{(*)}{(6月16日)其实,两年前他就开始露头了——见前一条脚注。1980年3月的那一幕,是Goulaouic--Schwartz讨论班;已引之注以及“掠夺的空白授权——或‘高尚行当’”一注(n\textsuperscript{o}~171$_{4}$,尤其打字稿第1088--1090页)都谈到它。}。反正,皮埃尔朋友深谙
此道,在一个人为的、由“行家们”斩钉截铁的共识拼凑出来的父亲身份的暧昧上,把戏做得
滴水不漏——而这早在那件新东西的分量得到公认之前就动工了。替身爸爸Kashiwara,1980年3月即已
登场,若不是早在六个月前的莱苏什讨论会上就已露头;十五个月后,在1981年6月那场难忘的
讨论会上,他被变走,不留痕迹(而且看来他自己也不太计较)。在这里,变走的手法炉火纯青:再引
进两位——这次不妨叫作“名义上的共同爸爸”——Beilinson与Bernstein;他们登台,纯粹是一道讲究
格式的保险条款——“喊停!”——因为当然,谁也不会想象孩子是他们俩谁做的(虽说两位都沾了
光……)。

与“Motif 运算”的类比,着实惊人!当人们还能拿出来的、关于 Motif 的“外行废话”的父亲身份,
知名度还太高(尤其七十年代初),无从下手时,Motif 瑜伽有两个\textbf{关键侧面},还不曾有一行
字公开发表过,哪怕以影射的方式。其一,“权的瑜伽”,早在1970年就被\textbf{超级爸爸}据为己有,
纹丝不动——反正被变走的那部分,不过是“猜想”,跟一句纯形式的影射差不多。另一侧面却毫无猜
测成分,六十年代下半叶就已打磨停当。按安排,由一个有点跟不上趟的无名学生,来陈述这套瑜伽
的启动机制——技术上不难,但(至少直到那位不受欢迎的生身爸爸“死亡”前后)这件活儿看来超
出了那可怜人的能力。

\pmark\origpage{1173}正是这位学生,Saavedra,是现成的替身爸爸:有幕后人暂时的担保撑腰,足以让一个一心
只想忘掉该忘之人的教团点头称是;但与此同时(妙就妙在这里),这位“爸爸”分明“没有分量”。
时候一到,谁也不会——Saavedra想必比任何人都更不会——起念头,宣称他可能是某种新哲学的父亲;只要
肯稍微想一想,这种假设简直滑稽……这里,功成身退的替身爸爸被请出去,要晚到十年:1982年,
难忘的 LN~900 出版。应当说,从1972年(第一场所谓“Motif 运算”中抬出“替身爸爸”)到1980年
(第四场“当值无名之辈”运算中同样应运而生的替身爸爸登场),多少水从桥下流过,再用不着绕
弯子了!耐人寻味的是,这里又插进来一位“名义上、装装门面的共同爸爸”,好让替代父亲身份(说
到底,就是木偶戏里的父亲身份……)与\textbf{真正的}父亲身份之间“平稳”过渡(而且没人显得在
抢镜头)。我确信,Milne对那双按他人意志摆布他的无形之手,看得不会比Beilinson与Bernstein更清
楚。人人都分到了碎屑;人人(至少是有发言权的人……)都有一切理由心满意足
\fnote{938}{(*)}{(6月16日)在他与Saavedra、与一位(暂时)在幕后的Deligne之间必然达成的“交易”里——这位Deligne,只等时候成熟便重新露面……——Saavedra的“份额”是:一张国家博士学位文凭,以及作为“数学讲义”这一显赫丛书作者而获得的相对知名度——这足以为他启程回国、远离他只远远瞥见过一眼的枯燥数学追求,铺好道路……}。

\pmark 这一切,昨夜还让我想起那第三场大运算——那场直接惠及“未来无所不在的爸爸”的运算:
“平展上同调”。此前我已使自己相信:这场运算的初始动机\fnote{939}{(*)}{关于这一点,见“那个公式”那组注(n\textsuperscript{o}s~169$_5$--169$_9$)。这最初的用意,后来大大扩展开来——尤其见“悼词(一)——或恭维”(n\textsuperscript{o}~104)与“宝石们”(n\textsuperscript{o}~170(iii))两条注。},是据为己有某一\textbf{不动点公式}的
打算:人们可以把某个父亲身份不受欢迎的“L 函数公式”呈现为该公式的一个平凡推论。麻烦在于,
所论的迹\dmark{}\origpage{1174}公式偏偏带着\textbf{同样}不受欢迎的父亲身份。幸好,还有另一位可能的父亲——一位好朋友(不点名,Verdier),他做
出了两份公式:一份太一般(但启发式地至关重要),另一份略显局促,却足以“罩住”人们想要的东
西。可朋友归朋友,这里的“父亲”既不能是朋友,也不能是那个不受欢迎的死者——因为事关“那
个”著名猜想的那把钥匙公式\fnote{940}{(*)}{当然是“那个”Weil猜想。关于这一点,见注“那个猜想”(n\textsuperscript{o}~169$_4$)。}。糟糕的是,L 函数公式与它那不合时宜的父亲身份,偏偏尽人皆知;所以
这里的难点不在朋友(朋友之间,最后总能商量……),而在死者。更倒霉的是,他对“推论”的证明,
1964年就白纸黑字发表在一场Bourbaki讨论班上——只是(万幸!)该公式的那桩例行情形(呃,
对不起,我是说关键情形!——这公式或迹公式,反正一码事,但这话千万说不得……\fnote{941}{(*)}{这两条公式,实际上互为直接推论。既然我对其中之一(L 函数公式)的父亲身份尽人皆知,Deligne便设法(在那份名叫“SGA~4½”的难忘文本里)把它呈现为另一条的推论,并且使尽浑身解数,以远比我这微薄的证明(以及随附的命题)艰难得多的魔术-行骗手法,把自己装扮成后者的父亲。这套戏法见已引的那组公式注(n\textsuperscript{o}s~169$_5$--169$_9$)——在我们这门可敬科学的年鉴里,这大概是绝无仅有的一场大变活人。}),当时还来不及被验证。

这里的手法,是拿那位朋友充当他的超一般公式的父亲(这原是分毫不差的事实——只除了一点:
他从来没费神去证明它……),同时悄悄把它与那累赘的死者所\textbf{显式}证明的公式(对那一条,自
然讳莫如深,一字不提)混作一团,再把那条超一般公式\textbf{拆毁}(说它是猜想的、不完全的、说
白了,没用的)。这叫溺毙一条鱼:叫读者压根不去想,翻一翻某套 SGA~5 讨论班(人们本来就责成
他“忘掉”它)会就这个问题说些什么。至于朋友那条(边上略局促、却完全有效的)显式公式,大家
一致同意也不再提了,只在一个文本的末尾留下一处暧昧而装门面的引用——那文本拖沓冗长、令
人生畏,世上没有哪个读者有勇气读到底。总而言之,可以说:“替身爸爸”(这里是Verdier)确实出
了场,但与其说是凭他对一项结果的默许的“父亲身份”——这里要\textbf{彻底变走}的,恰恰是\dmark{}\origpage{1175}死者的那
项结果——不如说是凭他的默许,参与了一场搅浑水、拆两条“孩子”的把戏:那两个孩子
确凿是他的,搅浑水是为了在混战中变走第三个孩子——那孩子的父亲不可宣之于口;如今是一个
孤儿,谁也找不回他,谁也不上心去找他\fnote{942}{(*)}{关于这一点,见注“变戏法者——或飞走的公式”(n\textsuperscript{o}~169$_8$)——以及“家庭相册”一注打字稿第1121页的脚注(**):它表明Deligne与Illusie这对好撒玛利亚人的变走-蒸发之计,是如何大功告成的。}。在这套手法里,Illusie扮演一个辅助角色,同方才
的“名义共同爸爸”相仿——只是他的父亲身份,同Verdier的一样,任何时候都不许落在那神圣的
\textbf{对 Frobenius} 的迹公式头上——那是唯一算数的,理当(当然,以应有的细腻)留给Deligne一
人;它落在那个要被变走的、不可宣之于口的孩子身上——Illusie以他特有的、堪称楷模的虔诚,参
与了这件事。

\subsubsection{e.\ Monsieur~Verdoux——或献殷勤的骑士}

\notehead{176$_{5}$}{}

可我还想回到Saavedra的“论文”上来。大约在我告别数学舞台之际——1970年初(如果我记得
不错)——Saavedra才装作终于真正“上钩”了:此前一两年,他一直显得三心二意。他对
我说,他已提炼出我最初给他出的那个命题的一种表述与证明,使它适用于基底环 $k$ 为\textbf{任意}环
的情形。他甚至给我画了一张证明的草图,我只半心半意地听了听——我那时几乎全部精力,都被我
正经历的人生转折占着。没有细核Saavedra的话,我有一种印象:他总算起步了,往后可以自食其
力。也许,在一个我对真正的科研指导已经几乎无暇顾及的时刻,我有些性急,把愿望当成了现实
\fnote{943}{(*)}{至少,同我离去之前的投入相比;但并不比我后来在大多数同事——那些担任科研指导的同事——身上看到的更少。}。此后,据我记忆,再没有他的音讯\fnote{944}{(*)}{我的记忆在此稍稍背叛了我——见注 n\textsuperscript{o}~176$_{7}$,其中对此有出乎意料的揭示。}。直到上星期,我还一直以为:我给他出的那个最低纲领,他总该完
成了,也许还稍微\dmark{}\origpage{1176}多做了些——把 Motif 的情形也处理了(照Deligne去年8月随其 Motif 带评注书目
写给我的信所言)。

\textbf{我这才刚刚发现:压根没有这回事}。这可怜人花了三四年功夫,竟在他称之为“塔纳卡范畴”
的东西的\textbf{定义}本身里(我指的是那个内蕴性质的定义),犯了一个粗糙已极的错误\fnote{945}{(*)}{错误的来由,是Saavedra心里对我所说的张量范畴的“基底环”有混淆:它绝不是使该范畴“线性”、使张量积“双线性”的随便什么环,而是典范的环 $\mathrm{End}(1)$(1 为范畴的单位对象)。我给Saavedra讲理论入门常识的时候,他想必“不在局中”到这般地步,这些话全从他头顶掠过,沉入遗忘。Deligne——他似乎多少接了我的班,当上了Saavedra的指导(显然另有一本账……)——却绝不肯替他纠正航向。这让他得以(十年之后)悄悄拆掉Saavedra的纸牌楼,以救世天使、以及(这一回又是)人人翘首以待的真爸爸的面目出现……}:而这个定义,正是要用来证明
它蕴含“Galois式”的、以一个合宜 gerbe 的表示写出的描述。他在引言里陈述的定理三(那篇引言,
按理至少应当\textbf{陈述}我交给他的那理论的四条根本定理)因此\textbf{平凡地为假}。Deligne与Milne以指出这个显赫的错误为乐事,提议把所论范畴的“新”定义取作 gerbe 式的描述(先天地明摆
着,这才是对的,必要时改一改内蕴描述就是……),并一本正经地追问:Saavedra的“定义”(把那
愚蠢的错误剔除之后)是否蕴含“他们的”(原文如此)\fnote{946}{(*)}{前引书第160页(我可不编瞎话!)。}——而这恰恰原本就是预定要构成Saavedra博士论文的那桩题目!

这一处境,是十足的于布爸爸的翻版!而且一上来就是三十六个花样。譬如,交给Saavedra的那
桩题目——唯一需要一点原创贡献的部分,哪怕多么微薄(为一个尽可能一般的基底环上的Galois-Poincaré范畴,提炼好的内蕴条件)——即使在(我相信我在遇到Saavedra之前很久就已处理过
\fnote{947}{(*)}{那是1964或1965年,即Saavedra那篇出名的“论文”(原文如此)之前七八年;比Deligne--Milne二人组赶来救驾——也\textbf{不做}这份微薄的活儿——还早十七八年:那是我对最微薄的学生所期待的唯一“原创”工作……})基底环\dmark{}\origpage{1177} $k=\mathrm{End}(1)$ 为\textbf{域}的情形
——也根本没有被处理。Saavedra的“论文”因此精确地在于:虔诚地照抄理论中(Grothendieck式瑜伽的启动之外)早已由我全部做完的那一部分——基底域上的那部分——并且,以一个
骗人的定义和一个假定理的“证明”顶替那先于一切的必做之事;那“证明”(Deligne以指出这一点为
己任——前引书第160页)归结起来,不过是一个简单的循环论证!

还不止这些。论文站不住——而答辩委员会毫无察觉!想必没有一个委员真正弄懂了所为何事。可
这并没有促使他们中任何一位向我示意:至少他们中该有一位,有能力为那项他们一本正经摆出架
势来评审的工作的严肃性,提供一份有效的担保\fnote{948}{(*)}{这支可悲的委员会的组成,终究要(在坚持读到那里的读者面前)在“第六根钉子”的终极之注 176$_{7}$ 里揭晓——在我棺材上的第六根钉子……}。答辩居然举行了,而且没有让我参与——这只能拜Deligne的担保所赐;
(如Saavedra的致谢所暗示的)在我实际上从舞台上消失之后,Deligne必定稍稍跟进过他的工作
\fnote{949}{(*)}{Deligne对一个默默无闻、想做论文的学生的这份突如其来的关切,是在该理论的自然父亲(且不受欢迎……)死后才出现的——何苦来哉;而那位学生(显然被任务压垮了……)本该陈述的,正是这门理论。}。

再说,打那时起,Deligne居然没有察觉这个错误,我简直无法想象——我深知他的敏锐与犀利,直至
最小的细节——而这绝非什么“小细节”!当然,我曾把已达成的瑜伽的全部精妙之处讲给他听;我
给他解释的最初几桩里,若不包含这个反例——他同Milne如今佯装当作最新发现拿出来的那个反
例,它早在我思考瑜伽的最初就为我所知(我最终打算把这瑜伽就叫“Grothendieck式”,不必去麻烦
Galois与Poincaré的大名……)——那是根本不可能的。他若听任“门生”(原文如此)的“论文”(再如
此)里留着这样一个粗糙的错误——只等时机一到,便可以纯然干脆地搞垮那位(暂时的)“替身爸
爸”——那想必不是没有充分理由的。昨天的思考,把那些理由已经摆得再明白不过。

\pmark\origpage{1178}也许有人会说,我在编故事:Saavedra所说的“帮助与建议”,并不意味着Deligne当真费神
稍稍认真读过引言里概括理论之实质的那四条命题\fnote{950}{(*)}{纤维函子上的滤套的那些结果除外——它们更技术性,更难压缩进一条醒目的命题。}。这些命题,他早在认识当事人之前就熟悉了;那么,那至多是一时轻率:为一项工作作
保,却连一刻钟的功夫都不肯花,验一验引言里预告的主要命题是否成立。可在我心里毫无疑义:Deligne确实费了这个神。因为这项工作\textbf{不是随便什么工作}——它出自一个多少有点找不着北、正
愁没论文的学生之手。继我之后(而且还在Serre之前),Deligne是最有资格感受那份形式主义的全部
分量的人:那是已故宗师留下的、未写出(至少未发表)的遗产的一个关键侧面。他尽管对这一侧面
摆出他惯常的满不在乎的架势\fnote{951}{(*)}{关于这些架势,以及它们所伺候的据为己有之术,见注“据为己有与轻蔑”(n\textsuperscript{o}~59$'$)。},骨子里他比谁都清楚所为何事。他,光辉的Deligne,极端的精英主义
者,居然肯费神跟进一个分明天资平庸的人的工作——这断然不是为了当事人漂亮的眼睛、为了帮
他拿到(按时下的共识,更按他所自诩奉行的、推到极致的严格标准)一纸\textbf{冒牌文凭}。

这话一出口,人们立刻面对一个奇怪的矛盾。一方面,一个据称多年全职投入本题目的人,犯了如
此显赫的错误,很难不把它解读为根本上无能的征兆——那个摆在他面前的问题,哪怕单看其技术
面(其实并不深奥),在答辩之时、在书出版之时,似乎根本没有被把握。另一方面,同一个学生,跟
我一两年没做出什么名堂之后,不到两年,忽然习得一门堪称令人叹服的数学素养:代数群结构理论,
一般域与\dmark{}\origpage{1179}实域上的、硬邦邦的概形理论、Hodge 理论、Motif……不止如此——我不记得读过任何
出自他手笔的数学文字,哪怕几页;而且我深知(对手段微薄的学生尤其如此)学会给数学
“成文”,绝非易事——可翻阅那本以他之名问世的书,它的“成色”之精良令我震动。我曾想,单就技
术而论,这份显然以与 EGA、SGA 同等的标准参考书自任的文本,本可以出自我的手笔,或Deligne
的手笔,或我那四五个才具出众、娴熟于把一套交织复杂的思想与事实精确、完整而优雅地陈述出
来的学生中任何一位的手笔。我很清楚,这样的行文功力,比数学素养更甚,绝非可以临时抱佛脚(除
非是天纵之才,如那位Deligne,及寥寥数人);它(倘若终究能获得)只有经年累月的练笔方能获得。
我自己就花了十年,而那时我同所要表达的实质的接触是极强的。这份接触,同Saavedra对他论
文题目的接触,当然不可同日而语——他写了关于这个题目的文字,却始终没有懂它;而这文字,后
来(至少到1982年……)竟然\textbf{是}“那一份”好的参考文献,为那套精微而关键的形式主义。分明
是两件干脆“对不上榫”的事……

昨夜起就掠过我心头的一个念头,如今,当我把处境白纸黑字对自己讲述一遍之后,带着明摆着的
力量回来了:Saavedra——我深知他的可能,尤其他的局限——不可能是那本光辉著作的作者;这
书以纯然技术的面貌(诚然,在这方面却又是穷尽的、针脚细密的),陈述一门完全超出他的“哲学”
的基础:也许头三章——其中两章主要是人人早已熟知的函子式泛论,第三章呈现的是书中中心观
念的、Saavedra版的十足的骗局——这几章,本该构成他从未完成的“最低纲领”;也许这几章全
然出自Saavedra的手笔。第三章尽管是十足的骗局,却仍足以\dmark{}\origpage{1180}让人窥见要到哪里去——即某些 $\otimes$-范畴的“Grothendieck式”眼光(姑且不点名),或“gerbe 式”眼光;它给后面的第四至
六章赋予意义。一旦接受了 gerbe 式的描述(在Deligne与Milne那双料海盗式的文本里,它被明智地
径直取作那些所谓“塔纳卡”范畴的\textbf{定义}),后三章便是要被据为己有的那套形式主义的
心脏。我猜,这几章是Deligne统统亲笔写的,或部分由他、部分由Berthelot写的;而且比我交给
Saavedra的笔记详尽得多,以致他几乎只消逐字照抄——如果曾有人要他费这番形式主义的神
的话。他准觉得“赢了”:人家“送”他一篇论文,外带一枚文凭——虽然他心里分明感到,他自己做的
那点东西(甚至是在自欺欺人地以为它站得住的情形下),对一篇国家博士论文来说,恐怕太单薄
了。而Deligne(又化装成撒玛利亚人……)赢了:这便是所需要的参考文献——纵然不为眼下,也为
“日后”(对会等待的人……),而且在那份参考文献里,那不受欢迎的名字,至少在一切实际的用处
上,不再出现。

为让欢乐登峰造极,我再补一句:那位名叫Saavedra的,似乎已从流通中消失,再无踪迹。去年,预
备寄发(我眼看就要寄出)《收获与播种》的印装本之时,我曾翻检《国际数学家人名录》——那册
子厚得很,“人人”都在(名录就是干这个的);唯独当事人不在:Saavedra之下没有,里瓦诺之下没
有,连Neantro之名下也没有(我凭良心翻过)。于是,这段历史有了黑色侦探小说的模样。想到那位
满面春风、和蔼可亲的Deligne——俨然又一位Monsieur~Verdoux(别名朗吕),以那本称心如意的“恰好的参考文献”遂
了他那番折磨人的心愿之后(比他朋友Verdier的还早四年!\fnote{952}{(*)}{关于它,见那道(理应如此)取名的注“好的参考文献”,n\textsuperscript{o}~82。})——人们不寒而栗:眼看着他让那桩恶魔般的阴谋的“证物”消失:那可怜的Neantro~Saavedra~Rivano本人,被放进奥尔马伊(Les~Ormails)一座精致的小壁炉里,长时间地焚化——那壁炉,
本就是特地为这类用途设计的\fnote{953}{(*)}{“奥尔马伊”(Les~Ormails)是法国高等科学研究所(IHÉS)住宅区的名字;我朋友皮埃尔——别名Monsieur~Verdoux,别名朗吕(还化装成献殷勤的骑士)——在那里及时接班,顶替了某位已故者:那人被排挤出这块地盘,又被我朋友所擅长的那种不动声色的暗算,抛进了暗无天日的地牢。住宅区有十来幢家庭小楼,和一栋大些的房子,带舒适的单人工作室;那些工作室想必不久也会各自配上一座小巧的全能壁炉……}。

\pmark\origpage{1181}我宽慰自己:没听说Kashiwara或Verdier也从这个世界消失——说真的,就在前天,我还拨通了后
者的电话,问他(没什么信心,似乎也没成功)能不能给我带来另一位“失踪者”的消息:人人谈论、
却看来无一人见过的——Jouanolou的论文。那篇论文,我仍然一无所知;但至少,Verdier似乎
还活着,尽管他本人就是一份“证物”——而且我怀着良好的希望:Neantro~Saavedra~Rivano也
一样。

\subsubsection{f.\ 肮脏的差事}

\notehead{176$_{6}$}{}

话虽如此,Saavedra论文史的那些于布式面向,我甚至还没有巡视完——看来那些不同寻常的论文与论文人,我是
收集上了瘾!到这里,我已有了揣测(不说“内心确信”吧):Deligne(由一位殷勤而义务的助手协助)
之所以装模作样,在Saavedra答辩十年之后一本正经“誊抄”其论文,想必只是“取回”他当年乐意
“借”给Saavedra一段时间的东西(这段时间,恰够Saavedra把论文过了、然后消失);所以这
不过是天理昭昭——只除了一点:他“借”来一段时间的,是从那位从不被点名的死者那里“借”来的。
可是,借死者的东西不还(那还了得!),风纪倒也无亏:这一头,也是万事大吉。

最妙的还在后头:即便在第二位前学生也过了这道手之后(而且还是我有过的学生中最光
辉的一位),我交给Saavedra的那个谦卑(微末)的问题——那曾是二十多年前我的出发点,
也是我自认当时就解决了的第一桩事,在所论 $\otimes$-范畴的定义环是域的情形——这个谦卑(微末)的问题,
直到今天,即便在那个情形里,居然仍然没有“被解决”!Deligne满足于点破Saavedra那个粗糙的错
误(那错误肯定十多年前就被看破了,他只是等待时机……)。他一边在128页上誊抄先前那本参考文
献\dmark{}\origpage{1182},却无心去修复这个错误。何必费这个神呢——追求的目的分明已经达到!要做到那一步,他
心里就必须在这场行动之中,除了据为己有的饥渴之外,另有别的东西:一份醒着的关切,一份对他
所处理的数学实质的\textbf{敬意},一种超出眼前“收益”的视野。

如果我在六四、六五年前后,曾费神为那些可以用“代数 gerbe”表述的 $\otimes$-范畴提炼一套“Grothendieck式”的瑜伽,而不满足于可用群概形描述的那些,那是因为在“打动”我最深的那个例子——一个域
上的 Motif——里,众所周知(凭Serre一个非常简单的论证):当该域特征 $p>0$ 时,不存在“有理
于 $\mathbb{Q}$”(甚至有理于 $\mathbb{R}$)的纤维函子。这在当时\textbf{逼着我的手},把理论表达在像 gerbe
与 lien 的形式主义这样“不正经”的东西里;同时当然,也要找到简单代数性质的内蕴判据,保证这种
“Galois式”或“Grothendieck式”的视野几乎“总是”行得通,而且花费极小。我所提炼的判据(如果
我没弄错,还给出了证明)——由 $k$ 的一个扩域 $k'$ 上纤维函子的存在给出——二十年后,文献里仍
然没有建立!直到今天,照Saavedra、Deligne之流写下的文字,纵然人们对(比方说)有限域上
“Motif 上同调类”的形式主义承认我们想承认的一切,半单 Motif 范畴(比方说)是“Grothendieck式”
的(或照这些先生的叫法,“塔纳卡的”)这一点,仍然没有建立(至少文献里没有)。前后共 418+128
=546 页的文字,出自Saavedra(由一位Deligne与一位Berthelot协助)继而Deligne与Milne之笔——
到头来却连二十年前我的出发点都没有抵达:那出发点,是让我自己相信,“Motif Galois群”\textbf{确
实存在}。

是啊,一个Deligne何必费这个神呢——他早已忘了那视野;他要的声望反正已到手;他用来做
\textbf{他的} Motif 理论(它当然同某位死者的理论毫无干系……)的那些域,全是特征零的域——于
是他那些所谓的“塔纳卡”范畴,统统是“中性”的(或“平凡”的)。照这个账,自然犯不上为 gerbe 之类
大做文章\dmark{}\origpage{1183}——那从此不过是障眼的烟幕。犯不上,除非是为了\textbf{据为己有一个}
\textbf{连灵魂与精神都已被遗忘的东西的字面}。而这场令人瞠目、可悲的故事的收梢,我看到,是这
样的:正如那套被埋葬了十五年的 Motif 视野的入门常识,还得由那个老朽来做——他刚巡礼完那
场光辉的葬礼及其绝技,就得自己把这份小活儿干了:他的学生们在他“死后”,还没有谁把它放在
心上。因为他们早就忙于扮演主子,忙得没有几天工夫,哪怕片刻,再当一回\textbf{仆人}
\fnote{954}{(*)}{我在这里有点性急,装作把我所有学生同他们当中最光辉的一位装进同一个口袋。我先向他们中那些不乐意与如此光辉的人士为伍的诸位致歉!无论如何,我高兴地记得Giraud:他硬着头皮接下了(不期而至的)读Contou-Carrère论文的活儿;那心境,对Contou-Carrère、对我,无疑是一种“服务”的心境,也许对数学共同体也是如此;关于这一点,见“耶稣与十二使徒”一注(n\textsuperscript{o}~19)最末一段,打字稿第151页。}。

\subsubsection{g.\ 一场屠杀中的五篇论文——或孝道}

\notehead{176$_{7}$}{}

(6月19日)到今天恰好两个月,我乘着一阵旋风,动手写前面那些注(4月19、20日的),带着那
个现成的名字“第六根钉子(钉进棺材)”(n\textsuperscript{o}s~176$_{1}$至176$_{6}$,不计本注——它也在其列)。此前一星期,Zoghman~Mebkhout刚把Saavedra那本书带给我——只消看上一眼,我就已经明白是怎么回事了。

我必须承认,这一发现带来的震动,几乎不亚于一年前(整整一年前)那本 Motif 掘墓“难忘文集”
(《数学讲义》第900种)带来的震动。更确切地说:去年的那份震动,某种意义上重新出现了,被这
个与那场掘墓紧密相联的“运算”的发现猝然重新点着;这场运算(一开始就明摆着)曾为它作了准
备,规模也相仿。我再一次被那种安详的无耻之感攫住(简直要说被噎得透不过气)——\textbf{同样的}无耻(这一点,凭种种不会
骗人的迹象,一开始也明摆着)——冲着一件与我血肉相系的东西,一件世上除我之外无人长久怀抱
与滋养的东西……那感受如此强烈,近乎焦虑,连我自己都吃惊。

\pmark\origpage{1184}自发的反应、自然的出口,本该像去年那样:趁感受还新鲜就说出来,由此径直进入这新的侧
面——那场由我的亲近者对我施行的活埋的新侧面。可我忍住了\fnote{955}{(*)}{我到底还是乘着当时的激动写了四五页,但九天之后(4月19日)写下正文时,那些痕迹已所剩无几。};因为Mebkhout的来访,需要我至少有一些闲暇;何况他
有些话要对我说——即便不那么触动我的要害,我感觉得到,它们对他而言同样是“要害”,对\emph{葬
礼}同样意味深长。再说,趁新知还热,把它们记下来,在我看来也重要;而那本出名的“葬礼之书”的
来龙去脉,即便我日后再来清理,也不会溜掉。所以,朋友刚走,我(4月15日至18日)就去写他的倒
霉遭遇,成如今“封神”末尾的那组注(n\textsuperscript{o}s~171$_{1}$至171$_{4}$)。

也就是说,在来到那出名的“第六根钉子”之前,我已经有时间恢复镇定。说实话,刚才重读最初几
页,在我对那桩新揭发的丑闻所作的挖苦(而稍带距离)的描述里,我找不出最初袭我的心绪的痕迹
——那心绪曾让我失眠一夜,而我当时多么需要睡眠。这一回,我是真切地感到了那份“过去的重量”!

那是6月10日,在给\emph{葬礼}画上那个出名的“句号”之后三天——葬礼随之重新开动!当然,我远
没有料到它开动到了什么地步——还有(差不多)三百页要写!写完“第六根钉子”的第六道注(“肮
脏的差事”)时,我满以为连“四种运算”也一并走完了——只除十来页(第三、四场运算的)要重抄
清稿、补上预定的脚注。再过几天,我想就可以把葬礼(III)的全部手稿送去打字了。

然而,随后几天里(也许就在我以为钉完最后一根“钉子”的次日或隔日),发生了一场出人意料的
剧情突变,对此我还得回头再谈。这里,我自发的冲动本该是立刻扑上去。我又等了两个月才动手,
当然不是不情愿\dmark{}\origpage{1185}——是有更急的事要备好送去打字。从头重读《四种运算》,发觉这里那里大有充实的必
要——后事如何,大家都知道了!

不管怎样,今天(只要不再有新的意外——摸摸木头!),终于到了那个黄道吉日:我给\emph{葬礼}
画上真正的句号——照实说了:写下最后几页的日子;它们理应构成我在《收获与播种》之中关于
葬礼的思考的尾声。此后,只还剩写那封充当《收获与播种》前言的“信”——然后我想休息几天,
受之无愧,而且也实在需要……

写完前面六道注几天之后,我得知了Saavedra答辩委员会的组成——那支我在“Monsieur~Verdoux——
或献殷勤的骑士”倒数第二道注里痛加讥讽(罪有应得)的委员会。论文于1972年2月25日在奥赛的
理学院答辩,委员会由\textbf{Demazure}(J.~Demazure,报告人)、\textbf{卡斯泰勒}(Castelle)与
\textbf{Grothendieck}组成。

要说“剧情突变”,这才是剧情突变!于布戏的加冕礼!而且,面对这份官方来源的信息,我好不容
易才相信——我连出席过这样一场答辩的一丝记忆的影子都没留下。Verdoux-朗吕先生的故事,愈发精彩了!我姑且给
Demazure打了个电话,问他可记得同我一起,给一个叫Saavedra的人当过论文答辩委员会。Demazure也记不得多少了,但好歹足以向我保证:答辩确凿举行过(何时如何,他自然说不太清),我们
两位都在场,还有卡斯泰勒(连这名字我都不记得了……)。他知道的也不多,只记得自己是论文报
告人。还是我告诉他:论文官方说来,是一份25页的文本(想必这让他的报告人差事轻松了不少)。
这一来,吃惊的反倒是他。他答应想法给我寄一份论文来。我本来很想知道它长什么样,可至今没有
等到——看来(照Demazure几个星期后终于对我说的)这份论文已无处可寻;也许他也没怎么费神
去找。反正,同我一样,他的文件里也毫无踪迹。不过,这只是个细节……

\pmark\origpage{1186}这一来,我可够瞧的了!我对那支分明无能的委员会——“一本正经地摆出架势评审一项他
们‘未必弄懂了所为何事’的工作”——所发的那些冷嘲热讽,如今都泼回我自己身上!可以想见,我多么想收
回那些讥讽,保住体面,撑住场面——可是不,那是作假。整场\emph{葬礼}里的作假已经够多了,犯不着
再添上我的一份。再说一遍:那些讥讽完全是有道理的。如今知道了委员会的组成,我甚至可以指明:
正是我,先于任何人,最应当受这些讥讽。毕竟,Demazure与卡斯泰勒当时记住的,主要是:这篇论文,
Saavedra是跟我准备的,或至少是跟我开的头,题目是我给的。按说“在局中”的是我,而他们信任
我。说不定,那出名的25页——Demazure据称做过报告人的——本来是站得住的;即便那桩显
赫的蠢事也在其中,它也不过是一份理论概述,不在局中、信任我的Demazure,绝无机会察觉。

至于我——除了上课,两年来实际已脱离数学——这场答辩,我草草了结,想必不过是在奥赛的两堂
课之间、某场“生存与生活”的集会、或某场关于附近(萨克莱)核废料的公开讨论之间,纯然一道
行政手续。确定无疑的是:两年来,Saavedra的工作我早已没有跟进,任何人的也没有;而且我对
Saavedra的工作站得住,毫不怀疑。这份确信从何而来,如今说不清了。同此前我所有的学生不同,
对Saavedra,我没有任何直接的凭据——凭他跟我一起已完成的工作——来判定他的严肃。难道那
些年,我竟把我的学术职责轻率到这般地步,凭他的面色,就信任了他?倘若那本书(同年出版,那
25页的论文想必是它的摘要)当时已经杀青,并让我形成了印象,那“看一眼”确实印象极佳,说
不定我连想都没想过要去核实:那本该构成Saavedra个人贡献的部分。也可能,甚至更可能(此事
我毫无记忆了),我信赖了Deligne的判断——我离去之后,是他跟进了这项工作
\fnote{956}{(*)}{我连Deligne曾跟进Saavedra的工作这件事,都没有记忆。这是4月里,看Saavedra那本书的引言时才知道的。}。

\pmark\origpage{1187}无论哪种情形,我都必须承认:我负有同样的责任——我把科学博士的头衔,授予了一篇二十
三年后显出是\textbf{冒牌货}的论文(借已引之注的说法)。然而,我自己在不知不觉中成了这场骗局的
工具、并因轻率作保而负有责任,这一事实,并不改变骗局之为骗局。它反倒显得愈发天才。因为说到
底,牵线者的真正动机,当然不是让一个落魄的论文新手廉价地得一顶头衔,然后改行退场——而是
让一个绝不糊涂的人,细腻地、不动声色地,把某种在我之内诞生、并且在他还没有听到(在数学
里)说出“gerbe”或“Motif”这类字眼之前,就被我孕育至成熟的那份视野的父亲身份,据为己有。正是趁
我骤然全力投入物种生存与其他种种最紧迫的美好事业之时(那位前学生兼朋友曾对我说,他不得
不同这些保持距离,因为他全身心绝对奉献给唯一的数学
\fnote{957}{(*)}{关于这一点,见注“兄弟与配偶——或双重签名”(n\textsuperscript{o}~134),尤其打字稿第614--615页。})——趁我的精力完全被别处占尽的时刻,我那位光辉的学生兼朋友,完成了这桩
真正独一无二的障眼大戏:让我成了我自己被褫夺的工具!我当时的心境,完全切断了对昔日数学兴
趣的牵挂,对我那些学生——Deligne首当其冲,他们自 SGA~5 讨论班结束起,就已开始按自己的路
数玩起了小把戏——报以盲目的信任:随便什么名字,只要安到那些我只剩遥远记忆的出名范畴上,
我都会说:是,阿门!正如我对Verdier宣布不会再有新式同调代数的书时说了是、阿门;正如Deligne向我宣布,我们合搞的 SGA~7 讨论班的一半,将突然更换父亲身份时,我说了是、阿门……然而,
受骗者乖乖地点头、毫无所觉地同意,并不改变诈骗的性质——只给它添上一重背信。而当Serre之
流的其他占卜师也从中各得其所、毫无保留地祝福其事
\fnote{958}{(*)}{这份再明白不过的祝福,见“家庭相册”一注部分 d.(“入殓——或自然的坡道”)。},这件事便有了非同寻常的维度\dmark{}\origpage{1188}——它成了一个圈子、一个时代的集体败坏——但它并不因此就光彩起来,无论多么天才;它的无耻,一分也不少。

正如剧情突变式的好事从不单行:得知我前学生Saavedra答辩委员会的组成才几天,关于Jouanolou论文的相应信息也到了——那也是一篇有点特别的论文,我在思考中这里那里谈到过它
\fnote{959}{(*)}{这篇论文,在注“团结”(n\textsuperscript{o}~85$^{1}$,打字稿第349页)的子注,以及“共同继承人……”(n\textsuperscript{o}~91,第387--388页)一注里谈到过。另见“学生与纲领”一节(n\textsuperscript{o}~25)。}。同Saavedra一样,他也从未想过给我寄一份他那出名的论文(“人人都引(自反常层会议以来),却无人见过”)——于是我最终给他写了封略嫌干巴的信(4月25日),就这
篇论文的离奇遭遇问了他几个问题。他几乎当即回了信(5月1日):对实质性的问题闪烁其词(因为
“重提往事总归非常痛苦”),但行政细节却精确之至:论文于1969年7月3日在法国高等科学研究
所(巴黎)答辩,委员会主席为\textbf{P.~Samuel},考试官为\textbf{迪斯米耶}(J.~Dixmier)、
\textbf{Grothendieck}与\textbf{Verdier}。我的通信人还添了一句,不无恶意:“据我所知,委员会全
体成员都到了!”(此事 J.L.~Verdier不久后也在电话里向我证实。)

\pmark 又是如此:对这场答辩,我连一丝记忆的影子都没有——它分明也是“偷偷”过的(不得不这
样糟蹋我自己的形象,抱歉!)\fnote{960}{(*)}{答辩的日期,我记得,应在我已经“脱离”数学、转而关心生物学(尤其分子生物学)之后不久。}。我以为答辩是在斯特拉斯堡(因而把它放在七十年代初——知道Jouanolou那几年在斯特
拉斯堡有职位),多半是因为Verdier在这篇论文上的一条隐晦的引文(1975年2月Bourbaki讨论班
讲义 n\textsuperscript{o}~464),引作“J.P.~Jouanolou,~Thèse,~Fac.~Sc.~Strasbourg”(无日期,无题目)。可他是答辩委
员会的一员,同我一样——他的记性也同我一样失灵,或者说得更准,任性,把举行答辩的 IHP\dmark{}\origpage{1189}(Henri~Poincaré研究所)安到了斯特拉斯堡?谁懂,谁就懂吧!

还是这位Verdier,好意把自己持有的那份论文寄给了我。起先,看着这包208张散页\fnote{961}{(*)}{索邦大学论文管理处存有一份215页的存档本——Verdier的这份看来缺了六页。说不定,存放在该管理处的那份,是世上唯一完整的一份——而且,人家还向我保证,是装订成册的。那里想必有一个专为弃儿论文装订的部门,它们都是散着零件到的……},我以为那是一份草
稿的复印件——我记得当年Jouanolou跟我做这篇迟迟不完的论文时,我亲手捧过并细细评注过的正
是它。可Verdier向我确认:这确凿是那篇倒霉论文的定稿本;它看来从未荣幸地印过三四份以上
(我那本,连同我的批注,想必回到了Jouanolou手里,从此再没见过……),也从未装订过。

后来,Jouanolou惠赐的稍详细的说明(6月3日的信),加上给Verdier的电话,让我重新入了一点
水。显然,Jouanolou对他的论文到达了一个“饱和点”:他从一开始就毫无信念地做(而我也懒得把处
境看清楚\fnote{962}{(*)}{见已引的“学生与纲领”一节,n\textsuperscript{o}~25。})。到1969年,他想必到了一个堵点,以致无力哪怕稍作修改,来照顾我的
许多批注。我只好认了,随它去。反正,重读一遍之后,我觉得这份文本代表了一份认真的、可用的
成形工作,虽远非完美——分明好于“聊胜于无”,在没有更合我意的文本的情况下,可以充当一份
必不可少的参考文献\fnote{963}{(*)}{不管怎样,它至今仍是世上唯一的、以导出范畴的版本陈述$\ell$-进系数理论的文本——而且偏偏无处可购,喜上加喜。电锯从那里经过了……}。

\pmark 当然,我做梦也想不到:Jouanolou会以他自己的方式,报复他跟我做这项工作时的三心二意
\dmark{}\origpage{1190}——亲手凿沉它,几乎抹掉那份我如此珍视的出名的“参考文献”的一切痕迹!这又是
一次“往事回头”,我实在没什么可抱怨的(尽管怨气并不是没有!)。在同Jouanolou的关系里,我要
的是在他身上找到“臂膀”,来推一辆规模可观的大车的轮子。我预先认定,他Jouanolou是我事业的参
与者,却从来没有停下来想一想,那些执拗的迹象分明告诉我:根本不是这么回事。诚然,是他自己选
择来跟我工作的(同一位名声显赫的“老板”共事,想必他觉得划算,却没料到自己陷进了什么……);
也是他自己,从我不吝扶持的一大片题目里,自由选定了他的工作题目(这些题目,当然全连着那辆
“大车”——而那车,骨子里多半引不起他半点兴致)。换句话说:同每个人一样,Jouanolou也在他自己
之内、在他自己的愿望与抉择的层面上,同某些矛盾搏斗——这里,是在他的工作里。

我自己的矛盾,不在我同工作的关系里,而在一种对任务的如此的偏执,以致我在学生身上只看得
见受欢迎的臂膀,无法想象他们中任何一位会在我与之共事的工作中心怀分裂。有了对\emph{葬礼}
的长考所给的这一份额外的回顾,我如今发觉:在我的学生里,“分裂”的远不止Jouanolou一个。但他
是一个极端的个案:他是其中唯一没能认同自己所选任务的人;他的工作是在没有信念、没有欢乐
之中完成的。我在这一处境中的责任,在于我不肯真正去知悉它,宁可把本该次要的东西(完成我的
任务)\textbf{摆在前头},而把要紧的东西(学生所“选”的任务真正\textbf{属于}他、并且怀着欢乐去
做)撂在后头。

这想必就是为什么,Jouanolou也是我唯一曾在其身上察觉到怨怼的前学生(当然,那怨怼从不报姓
名)。怀怨是一个出口、一剂解药;它当然一无用处,只除了回避自己的问题(而人们很少往深里追
究)。可它毕竟是事出有因的;今天(二十年之后)我收割到它的某些果实,无权抱怨。

\pmark\origpage{1191}不到两个月里,接连面对Saavedra论文、继而Jouanolou论文的这些不寻常的插曲,让我切
身感到一件在《收获与播种》第一部分里才刚刚瞥见的事:早在我离去之前,以及紧接着的那些年
里,我的学生们和我之间,并不像我想当然的那样一切安好。譬如,在跟我做到国家博士论文层级的
学生们所通过的十二篇论文中,\textbf{四}篇赫然是针对宗师的“葬礼论文”!它们相继出现在1967至
1972年的五年之间,其中两场“葬礼论文”在我离去之前就已举行。第一场是Verdier1967年的论文,
缩减成一份28页的摘要,是我引入的新同调代数被埋葬的前奏;那套代数,Verdier曾受命去发展。
此事已有相当周详的交代\fnote{964}{(*)}{关于这一点,尤其见“赊账的学位论文与全险”与“随心所欲的荣耀——或暧昧”两条注(n\textsuperscript{o}s~81、170(ii))。},不必重复。第二场是Jouanolou1969年的论文,它为$\ell$-进上同调的形式主义——从导出范畴
与三角范畴的观点看(这对六种运算分明至关紧要;Verdier本该提供基本参考)——的埋葬盖了
章。第三场是Deligne1970年的(?)论文:一篇十足光辉的论文,也深深扎根于他从我这里得来的思
想\fnote{965}{(*)}{这是Deligne的“Hodge 理论 II”。关于这项工作在 Motif 瑜伽、在我的“系数理论”视野(包括一种“霍奇系数”的理论)之中的扎根,我在注“把一个个 i 点清”(n\textsuperscript{o}~164,尤其打字稿第739--740页)以及子注 n\textsuperscript{o}~164$_{1}$(第805--806页)里给出细节。同 M.~Raynaud与 C.~Contou-Carrère一样,Deligne自选了他的工作题目——包括论文题目——并没有等我给他出题,而且完全独立地做了这项工作,甚至在几乎完成之前都没跟我提起。尽管如此,他的工作(论混合 Hodge 结构)扎根于我的思想,比Raynaud与Contou-Carrère更深:后两位主要使用我带来的语言与技术,而他们各自追究的问题域,则是全然原创的。
的确(照今天刮的风),思想随风而去——尤其当它们没有发表的话(几天前,Serre又斩钉截铁地向
我解释过一遍……)。}——而我的名字在文中连提也不提!第四场是Saavedra的论文——方才长
篇谈过:不是那位推定的作者\fnote{966}{(*)}{这至少是我在“Monsieur~Verdoux——或献殷勤的骑士”(n\textsuperscript{o}~176$_{5}$)那一注里达到的确信。},而是别人,以他为人熟知的娴熟手法,陈述了第三个人关于 Motif Galois群的思想与\dmark{}\origpage{1192}结果(经由一套完整的所谓“塔纳卡”范畴理论——而且是整整四种!),却对我这卑不足道、亡故之人只字不提!

这四场葬礼式论文(它们正是“四种运算”——大写的那种!——的序曲)分明以许多方式互相勾连
\fnote{967}{(*)}{往深里探测这些联系,当然有趣——但如我在几行之下所说,现在不是时候。}。它们相续于不到五年之间,始于紧接 SGA~5 讨论班结束的那一年。那场讨论班,看来确是
我前学生们身上掘墓心性的起点与集结之地——而且早在我离去之前!这些心性先于我的离去而存
在,这是一个耐人寻味的情况,关乎\emph{葬礼}的那个“第二平面”——由“从前的”全体前学生组
成;这个情况,我还没有真正整合进一个总体的理解。眼下,这“第二平面”似乎是三个平面中最不被
理解的一个。但现在不是重新开启思考的时候;接下来的几个月,想必不会缺少新的材料,它们会由
我的前学生们亲手送来。到那时,再把它们拼成“第二平面”的一幅活的总体图景。

还有第五篇论文\fnote{968}{(*)}{在共十四篇论文之中——十四位(无论“从前的”还是“后来的”)跟我做到国家博士论文层级的学生。也就是说,按数目算,\textbf{三篇之中就有一篇以上}是葬礼论文——已经不算糟了!},对我而言,它也插进葬礼论文的序列,不过是一篇“之后”的论文——比前一序列晚了十
年。那是Contou-Carrère的论文,1982年12月通过;它也在不止一个方面很特别。它不同于前四篇之处在
于:Contou-Carrère为讨好要人们、为求得原谅自己曾多多少少当过我的学生而作出的英勇的掘墓努
力,并没有能避免Verdier(他以为选他做论文导师是聪明的
\fnote{969}{(*)}{而且那时候,我还一直(照Contou-Carrère本人向我保证的)以为自己是他的官方论文导师。我得知还有一位“平行”导师的存在——在这一对里,本该是我充当“备位”导师,以防万一……——是在Contou-Carrère不得不回头找我的时候;与此同时(局势已经变得有点太那个了),他向我揭示了Verdier扮演的角色。经年累月地玩着这种没谱的把戏,Contou-Carrère最后几乎不再做数学,难怪。可以说,他不是唯一的一个……})突然装模作样地要把他\textbf{搞沉}\dmark{}\origpage{1193},不宣而战——于是,不得已,他又回头找我。我本不必
出场充当论文导师:Contou-Carrère自己找到了题目、自己发展了方法,我也没有跟进他的工作,何况
他的工作所在的语境(约化群概形)我已有些生疏。尽管如此,他工作的初始想法——为 Schubert 闭
链的闭包搞一种“等变”的奇点消解方法——直接受启发于我(约1975或76年)向他原原本本解释过的
一个想法:关于约化群在自身上的伴随表示,其轨道闭包的典范的、同步的奇点消解
\fnote{970}{(*)}{六十年代末,我曾被Brieskorn关于所谓“有理”(曲面)奇点的漂亮工作、及其与某些单根系(根全等长的那些)的联系所吸引;我给自己出了一个问题(不用说,有点想入非非):用与根图相应的单代数群,直接描述一个有理奇点。就这样,我得到了所论奇点消解的一份非常简单(说白了,甚至是显然)的几何描述,靠的是 Killing 对;外加一整套精彩的猜想作为钥匙,后来我有些忘了,当年讲给过凡愿听的人。可我什么也没发表;而照Serre刚向我惠赐的新公理——归第一个捡到的人所有——我发觉,有人捡了不少,难免。有时候,换换公理,真方便……}。不用说,Contou-Carrère早就嗅准了上流社会里风往哪儿吹——那上流社会,他正合法地盼着跻
身;对这层渊源,他守口如瓶。倘若大家又都开始提这类不可捉摸的东西,那还了得:一个\textbf{观念}
(还没发表呢),据说是为了\textbf{激发}另一个(您给评评理……)——当然,除非被光荣引用的,正
是那些让所呈工作增光的名字之一(到那时,自然也无须解释为什么致谢滔滔,那致谢想必是名正
言顺的……)。

\asterisms

\centerline{\headfont “四种运算(对一具尸体)”终}

